\documentclass[11pt]{article}
\usepackage[a4paper,margin=18mm]{geometry}
\usepackage[no-math]{fontspec}
\setmainfont{Latin Modern Roman}
\usepackage{amsmath,amssymb,amsthm,xfrac,xspace,hyperref,graphicx,comment}
\excludecomment{DefTralics}
\providecommand{\dpl}{\displaystyle}
\providecommand{\backmatter}{}
\providecommand{\loccit}{\emph{loc. cit.}}

\providecommand{\bibliofont}{\small}
\newcommand{\ie}{i.e.,\xspace}
\newcommand{\eg}{e.g.,\xspace}
\newcommand{\etc}{etc.\xspace}
\newcommand{\cf}{cf.\xspace}
\newcommand{\resp}{resp.\xspace}
\newcommand{\smaller}{\small}
\newcommand{\Subsection}{\subsection}
\newcommand{\Subsubsection}{\subsubsection}
\newcommand{\skpt}{\smallskip}
\emergencystretch=3em
\numberwithin{equation}{section}
% Begin literal retained input author-preamble.tex
\usepackage[scr=rsfs,cal=euler]{mathalfa}
\makeatletter
\def\p@enumii{}
\makeatother
\hyphenation{semi-stable ortho-g-o-nal togeth-er}
\newcommand{\Fil}{\textit{Fil}}
\newcommand\hto{\mathrel{\lhook\joinrel\to}}
\newcommand{\onto}{\mathchoice{\to\hspace*{-5mm}\to}{\to\hspace*{-2.85mm}\to}{\to\hspace*{-3mm}\to}{\to\hspace*{-3mm}\to}}
\newcommand\To[1]{\mathchoice{\xrightarrow{\textstyle\kern4pt#1\kern3pt}}{\stackrel{#1}{\to}}{}{}}
\newcommand{\sbullet}{{\scriptscriptstyle\bullet}}
\newcommand{\Psfrac}[2]{(\sfrac{#1}{#2})}
\newcommand{\spfrac}[2]{\sfrac{#1}{(#2)}}

\newcommand{\lconnection}{$\lambda$\nobreakdash-con\-nec\-tion\xspace}
\newcommand{\lconnections}{$\lambda$\nobreakdash-con\-nec\-tions\xspace}

\newcommand{\grsemistable}{gr\nobreakdash-semi\-stable\xspace}
\newcommand{\Dol}{\mathrm{Dol}}
\newcommand{\dR}{\mathrm{dR}}
\newcommand{\Fib}{\mathrm{Fib}}
\newcommand{\HN}{\mathrm{HN}}
\newcommand{\Hod}{\mathrm{Hod}}
\newcommand{\Rees}{\mathrm{Rees}}

\usepackage{tikz-cd}

\DeclareMathOperator{\codim}{codim}
\DeclareMathOperator{\Ext}{Ext}
\DeclareMathOperator{\Filt}{Fil}
\DeclareMathOperator{\Hom}{Hom}
\DeclareMathOperator{\GL}{GL}
\DeclareMathOperator{\Rep}{Rep}
\DeclareMathOperator{\SL}{SL}
\DeclareMathOperator{\Spec}{Spec}
\DeclareMathOperator{\Vect}{Vect}

\newcommand{\sO}{{\mathcal O}}
\newcommand{\PP}{{\mathbb P}}
\newcommand{\rk}{\mathrm{rk}}
\newcommand{\Gr}{\mathrm{Gr}}

\theoremstyle{plain}
\newtheorem{thm}{Theorem}[section]
\newtheorem{cor}[thm]{Corollary}
\newtheorem{lem}[thm]{Lemma}
\newtheorem{prop}[thm]{Proposition}
\newtheorem{ques}[thm]{Question}
\theoremstyle{definition}
\newtheorem{defn}[thm]{Definition}
\newtheorem{rem}[thm]{Remark}
\newtheorem{examp}[thm]{Example}
\newtheorem{conj}[thm]{Conjecture}
\newtheorem{notn}[thm]{Notation}

% End literal retained input author-preamble.tex

\begin{document}
\begin{center}{\Large Stable item 2281}\end{center}
\paragraph{Source and attribution.}
Mao Sheng, Hao Sun, Jianping Wang, \emph{On the existence of gr-semistable filtrations of orthogonal/symplectic $\lambda $-connections}, Journal de l’École polytechnique — Mathématiques 11 (2024), publisher version, DOI 10.5802/jep.276. \href{https://jep.centre-mersenne.org/articles/10.5802/jep.276/}{Publisher article}. Authors' text: \href{https://creativecommons.org/licenses/by/4.0/}{CC BY 4.0}.
\paragraph{Editorial problem boundary.}
Construct only the characteristic-zero rank-five orthogonal connection in Example5.7 below, under the original smooth complex projective curve, stable rank-two bundle of positive degree and genus bound. Prove that its orthogonal connection exists, is semistable and admits no gr-semistable orthogonal filtration, with the exact nonsplit tensor-product construction and isotropic destabilizer conventions. Retain the full new quasi-filtration obstruction and the odd-rank equivalence used to pass from quasi to graded semistability. No general rank classification or positive-characteristic strengthening is selected.
\paragraph{Difficulty (editorial): H3.}
The cited orthogonal-connection criterion alone does not produce a connection on the chosen nonsplit tensor product: the authors exclude every positive-degree isotropic decomposition, using stability and a splitting contradiction for the rank-two case. They then identify its stable rank-two destabilizer and show that connection invariance of its perpendicular would force a forbidden degree-zero connection. The new quasi-filtration obstruction and odd-rank filtration equivalence transfer that obstruction to every gr-semistable orthogonal filtration, beyond the allowed bundle and connection foundations.
\paragraph{Editorial transcription note.}
Original author language, mathematics and ordering are preserved. Publisher front matter is replaced by this independent article wrapper. Bibliography is recovered from matching cached publisher HTML, including its TeX formulas. No new proof or mathematical repair is supplied. Surrounding original results are retained for conservative context and are not extra selected corpus claims.
\clearpage

% Begin literal retained input author-body.tex
\section{Introduction}

Let $X$ be a smooth projective variety over an algebraically closed field $k$. Given $\lambda \in k$, a \lconnection on a vector bundle $V$ over $X$ is a $k$-linear morphism
\[
\nabla: V \to V \otimes \Omega_X,
\]
satisfying the so-called $\lambda$-Leibniz rule. It is integrable if it satisfies the integrability condition $\nabla^2=0$. When $\lambda=1$, $(V,\nabla)$ is a flat bundle, while when $\lambda=0$, the pair is a Higgs bundle. When $k=\mathbb{C}$, Simpson considered the moduli space $\mathcal{M}_\Hod(X,r)$ of (Gieseker) semistable integrable \lconnections on $X$ of rank $r$, and showed that there is a natural morphism
\begin{align*}
\mathcal{M}_\Hod(X,r) \to \mathbb{A}^1, \quad (V,\nabla,\lambda) \mapsto \lambda
\end{align*}
such that the preimage of $\lambda=0$ is the Dolbeault moduli space $\mathcal{M}_\Dol(X,r)$ and the preimage of $\lambda=1$ is the de Rham moduli space $\mathcal{M}_\dR(X,r)$ \cite[Prop.\,4.1]{Si97}. There is a natural $\mathbb{G}_m$-action on $\mathcal{M}_\Hod(X,r)$, \ie
\begin{align*}
t \cdot (V,\nabla,\lambda) := (V, t\nabla, t\lambda),
\end{align*}
and consider the limit $\lim_{t \to 0} (V, t\nabla, t\lambda)$. In~the case of curves (hence Gieseker semistability coincides with slope semistability), the limit $\lim_{t \to 0} (V, t\nabla, t\lambda)$ always exists \cite[Lem.\,4.1]{S10}. Moreover, Simpson gave an effective approach to find such a limit, which is called \emph{iterated destabilizing modifications} \cite{S10}. We~briefly review the background of this approach. Let $(V,\nabla)$ be a flat bundle. An effective finite decreasing filtration of subbundles
\begin{align*}
\Filt: 0=\Fil_m \subset \Fil_{m-1} \subset \cdots \subset \Fil_0=V
\end{align*}
is called a \emph{Hodge filtration} of $(V,\nabla)$ if $\Filt$ satisfies the \emph{Griffiths transversality condition}:
\[
\nabla(\Fil_{i+1})\subset \Fil_{i}\otimes \Omega_X, \ i \geq 1.
\]
From a Hodge filtration we shall obtain a well-defined Higgs bundle $(E,\theta)$ by taking grading, where $E: = \Gr_{\Filt}(V)$ and $\theta$ is induced by $\nabla$. If $(E,\theta)$ is semistable as a Higgs bundle, then we say the Hodge filtration is \emph{\grsemistable}. In~fact, Simpson proved that \grsemistable filtrations always exist, which gives an alternative way to prove that the limit $\lim_{t \to 0} (V, t\nabla, t\lambda)$ exists. The notion of a \grsemistable Hodge filtration (with respect to a chosen polarization) has been generalized to a high dimensional base (\cite{LSZ19,L13}). Here we shall make the existence of a \grsemistable Hodge filtration as a defining property of a \lconnection, which is termed as \emph{gr-semistability} (see Definition \ref{gr-semistable}). It has the following fundamental significance.

\begin{thm}[{Simpson \cite[Th.\,2.5]{S10}, Lan-Sheng-Yang-Zuo \cite[Th.\,A.4]{LSZ19}, Langer \cite[Th.\,5.5]{L13}}]\label{the case of flat bundles}
Let $k$ be an algebraically closed field and $X$ a smooth projective variety over~$k$, equipped with an ample line bundle $L$. Let $(V,\nabla)$ be an integrable \lconnection over~$X$. Then $(V,\nabla)$ is $\mu_L$-semistable if and only if it is $\mu_L$-\grsemistable.
\end{thm}
The above theorem provides an alternative way to understand semistability for flat bundles in positive characteristic and the result plays a crucial role in the theory of semistable Higgs-de Rham flows \cite{LSZ19} and then its subsequent applications. Furthermore, when $L$ is clear in the context, we~shall omit $\mu_L$ when we are speaking about slope semistability. In~fact, we~shall simply speak of semistability as we shall only deal with slope semistability throughout the paper.

In our attempt to generalize the theory of semistable Higgs-de Rham flows to the case of principal bundles, we~found a different phenomenon, as we shall explain below. On the other hand, our study has also been driven by the following question of Simpson, which he proposed after he proved the existence of the limit $\lim_{t \to 0} (V, t\nabla, t\lambda)$ in the case of principal bundles (\cite[Cor.\,8.3]{S10}).
\begin{ques}[{\cite[Quest.\,8.4]{S10}}]\label{quest_simp}
How to give an explicit description of the limiting points in terms of Griffiths-transverse parabolic reductions in the case of principal bundles?
\end{ques}
In this paper, we~shall study this question, together with his proposal, in the special case of orthogonal/symplectic bundles.

During the study of this problem, we~find the following example indicating that the approach of iterated destabilizing modifications does not work any more in the case of orthogonal bundles, and thus principal bundles.

\begin{examp}[Example \ref{exmp_char_zero}]
We find an example of an orthogonal connection of rank $5$ in characteristic zero, which is not \grsemistable. We~refer the reader to Example \ref{exmp_char_zero} for more details.

Let $X$ be a complex smooth projective curve with genus $g \geq 2$. We~fix a rank $2$ stable vector bundle $M$ on $X$ with $\deg(M)=d>0$, and we suppose $g > d+1$. With respect to the above setup, we~can choose a nontrivial extension
\[
0 \to \mathcal{O}_X \to M' \to \det(M^{\vee}) \to 0.
\]
Then we define $E := M \otimes M'$ together with a natural orthogonal structure given as follows
\begin{multline*}
\langle \,, \, \rangle_E : E \otimes E \To{\cong} M \otimes M \otimes M' \otimes M' \to \wedge^2 M \otimes \wedge^2 M' \\
\To{\cong} \det(M) \otimes \det(M^{\vee}) \cong \mathcal{O}_X,
\end{multline*}
and then we obtain a rank $4$ orthogonal bundle $(E, \langle \,, \, \rangle_E)$. We~construct an orthogonal connection $\nabla_0$ on $E$ and define a rank $5$ orthogonal connection $(V, \langle \,, \, \rangle, \nabla)$ as follows:
\begin{itemize}
\item $V:= E \oplus \mathcal{O}_X$,
\item the orthogonal structure $\langle \,, \, \rangle$ on $V$ is induced by $\langle \,, \, \rangle_E$ and the trivial one on $\mathcal{O}_X$,
\item $\nabla:=\nabla_0 \oplus \mathrm{d}_\mathrm{can}$, where $\mathrm{d}_\mathrm{can}: \mathcal{O}_X \to \Omega_X$ is the exterior differential.
\end{itemize}
This orthogonal connection $(V, \langle \,, \, \rangle, \nabla)$ does not admit \emph{any} orthogonal Hodge filtration so that its associated graded orthogonal Higgs bundle is semistable.
\end{examp}

The example shows that the analogue of Theorem \ref{the case of flat bundles} is \emph{false} in the case of orthogonal integrable connections. That is, the notion of gr-semistability is strictly stronger than that of semistability for orthogonal integrable connections. Moreover, Alper, Halpern-Leistner and Heinloth introduces a powerful semistable reduction theorem for $\Theta$-stability condition \cite[Th.\,6.5]{AHLH23}, which can be used to prove the existence of such a limit for the corresponding moduli stacks under a certain base change. Based on Alper--Halpern-Leistner--Heinloth's result, our example actually shows that if we do not allow to choose a base change, such a limit may not exist. We~will discuss this property in a future project about Higgs-de Rham flow for principal bundles. To~the positive side of this study, we~obtained the following theorem, which is the main result of the paper.
\begin{thm}[Theorem \ref{thm_gr_and_quasi_gr_odd} and \ref{thm_even}]
Let $k,X,L$ be as in Theorem \ref{the case of flat bundles}. Let $D$ be a reduced normal crossing divisor on $X$. Let $(V,\langle \,, \, \rangle,\nabla)$ be a logarithmic orthogonal/symplectic \lconnection. Then it is \grsemistable if and only if it is semistable and quasi \grsemistable.
\end{thm}
\begin{rem}
We also consider $\Omega_X(\log D)$-valued \lconnections, which is termed as logarithmic \lconnections. Note that the integrability condition on a \lconnection is not required in the statement.
\end{rem}

The definition of quasi gr-semistability is somewhat technical. See Definitions \ref{defn_gr_semi_odd} and \ref{defn_gr_semi_even} for a precise formulation. Roughly speaking, the quasi gr-semistability is a necessary and sufficient condition, under which the method of iterated destabilizing modification works.

The structure of the paper is manifest: after laying down the necessary setup and terminologies in Section~\ref{sect_pre}, we~establish the main result in Section~\ref{sect_odd} and Section~\ref{sect_even}. In~Section~\ref{sect_small_rank}, we~explicate the quasi gr-semistability in small ranks, viz. $\rk(V)\leq 6$.

\subsubsection*{Acknowledgements}
The authors would like to thank the anonymous referee for numerous helpful comments and suggestions, which improved the quality of the present article.

\section{Preliminaries}\label{sect_pre}

Let $X$ be a smooth projective variety over an algebraically closed field $k$ with a fixed reduced normal crossing divisor $D=\sum_{i=1}^{k}D_i$. Denote by $\Omega_X$ the cotangent sheaf of~$X$. In~Section~\ref{subsect_rees}, we~review the correspondence between Rees $G$-bundles and filtered $G$-bundles studied in \cite{Lo17a,Lo17b}. The correspondence gives a way to study filtered $G$-bundles by taking an appropriate faithful representation $G \to \GL(W)$ and consider the corresponding filtered bundles. Based on this idea, we~give the definition of filtered orthogonal/symplectic sheaves, which is also considered in \hbox{\cite[\S5]{GS03}}. Then, in Section~\ref{subsect_log_lambda_conn}, we~introduce the main objects studied in this paper, the \emph{Hodge filtered logarithmic orthogonal/symplectic \lconnections} $(V, \langle \,, \, \rangle, \nabla, \Filt)$ (Definition~\ref{defn_na_Hod_fil}), where $(V, \langle \,, \, \rangle)$ is an orthogonal/symplectic sheaf, $\nabla$~is an orthogonal/symplectic \lconnection and $\Filt$ is an orthogonal/symplectic filtration obeying Griffiths transversality. In~the end of this section, we~discuss the stability condition of $(V, \langle \,, \, \rangle, \nabla)$ and its Harder--Narasimhan filtration.

\subsection{Rees construction}\label{subsect_rees}
In this subsection, we~briefly review Rees construction, which gives a correspondence between filtered $G$-bundles on $X$ and $\mathbb{G}_m$-equivariant $G$-bundles on $X \times \mathbb{A}^1$. We~refer the reader to \cite{L13,Lo17a,Lo17b} for more details.

Let $V$ be a vector bundle (locally free sheaf) of finite rank on $X$ together with a filtration
\begin{align*}
\Filt: 0 = \Fil_m \subseteq \Fil_{m-1} \subseteq \dots \subseteq \Fil_0 = V
\end{align*}
of subbundles on $X$. Furthermore, we~always assume that $\Fil_i =0$ for $i \geq m$ and $\Fil_i = V$ for $i \leq 0$. Such a pair $(V,\Filt)$ is called a \emph{filtered bundle} on $X$.

Now we consider Rees bundles. A \emph{Rees bundle} on $X$ is a $\mathbb{G}_m$-equivariant bundle~$\mathcal{V}$ of finite rank on $X \times \mathbb{A}^1$ (with respect to the natural $\mathbb{G}_m$-action on $\mathbb{A}^1$). Let $\mathbb{A}^1 = \Spec \, k[t]$. As proved in \cite[Lem.\,2.1.3]{Lo17b}, a Rees bundle on $X$ is equivalent to a triple $(V,\mathcal{V},\varphi)$ such that
\begin{itemize}
\item $\mathcal{V}$ is a vector bundle over $X \times \mathbb{A}^1$,
\item $V$ is a vector bundle over $X$,
\item $\varphi: \mathcal{V}|_{t \neq 0} \To{\cong} \pi_X^* V$ is an isomorphism, where $\pi_X$ is the projection to $X$, such that there exist generators $\widetilde{v}_i$ for $\mathcal{V}$ and $v_i$ for $V$ satisfying $\varphi(\widetilde{v}_i) = v_i \otimes t^{n_i}$, where $n_i \in \mathbb{Z}$.
\end{itemize}
Applying the argument of \cite[Prop.\,2.1.5]{Lo17b}, we~have the following equivalence.

\begin{lem}\label{lem_rees_fil}
The category $\Filt_X$ of filtered bundles on $X$ is equivalent to the category $\Rees_X$ of Rees bundles on $X$. Moreover, the equivalence holds as tensor categories.
\end{lem}

\begin{proof}
We only give the correspondence briefly. Given a filtered bundle $(V, \Filt)$, we~define a bundle
\begin{align*}
\mathcal{V}: = \sum t^{-i} \Fil_{i} \otimes_{\mathcal{O}_X} \mathcal{O}_{X \times \mathbb{A}^1}
\end{align*}
on $X \times \mathbb{A}^1$ together with a natural $\mathbb{G}_m$-equivariant structure induced from that on~$\mathbb{A}^1$. Clearly, $\mathcal{V}$ is a Rees bundle on $X$.

On the other direction, given a Rees bundle $\mathcal{V}$ on $X$, denote by $V: = \mathcal{V}|_{t = 1}$ the restriction to $t=1$. We~will give a filtration on $V$ by working locally on an affine open subset $U \subseteq X$, over which $V(U)$ is a free $\mathcal{O}_X$-module. Let $\{v_i\}$ be a set of generators of $V(U)$ and let $\{\widetilde{v}_i\}$ be a set of generators of $\mathcal{V}|_{t \neq 0} (U \times \mathbb{G}_m)$ such that $\varphi(\widetilde{v}_i) = v_i \otimes t^{n_i}$ under the isomorphism $\varphi: \mathcal{V}|_{t \neq 0} \cong \pi^*_X V$. Since $V(U)$ is free, we~can assume that $\{v_i\}$ is a basis of $V(U)$. Then, given an integer $i$, we~define a free submodule $\Fil_i(U) \subseteq V(U)$, which is spanned by $v_j$ such that $n_j \leq -i$. In~other words, $\Fil_i(U)$ includes all sections $v \in V(U)$ such that the vanishing order of $\widetilde{v}$ at $t=0$ is at least $i$. Clearly, $\Fil_{i+1}(U) \subseteq \Fil_{i}(U)$ and we obtain a filtration of $V(U)$. By patching together $\Fil_i(U)$, we~obtain a subbundle $\Fil_i \subseteq V$, and thus a filtration of subbundles of $V$ as desired.
\end{proof}

\begin{rem}\label{rem_graded}
Given a filtered bundle $(V, \Filt)$, we~can define a graded bundle $\Gr_{\Filt}(V) := \bigoplus_{i=0}^{m-1} \sfrac{\Fil_i}{\Fil_{i+1}}$, which is also the restriction of the corresponding Rees bundle $\mathcal{V}$ to $t= 0$.
\end{rem}

Now we consider the case of principal bundles. Let $G$ be a connected reductive group over $k$. Similar to the case of vector bundles, a \emph{Rees $G$-bundle} is defined as a $\mathbb{G}_m$-equivariant $G$-bundle on $X \times \mathbb{A}^1$. Now we will give an equivalent definition of Rees $G$-bundle in the viewpoint of Tannakian categories. Denote by $\Rep(G)$ the category of $G$-representations with values in free $k$-modules of finite rank, which is a (rigid) tensor category. Recall that a $G$-bundle is equivalent to a faithful exact tensor functor $\Rep(G) \to \Vect_X$ (see \cite[\S 6]{Si92} for instance). Now we choose a representation $\rho : G \to \GL(W)$ and a Rees $G$-bundle $\mathcal{E}$, and define the associated bundle $\mathcal{E} \times_G W$, which is a Rees bundle on $X$. This construction induces a fiber functor $\Rep(G) \to \Rees_X$. Lovering proved the following equivalence of categories.

\begin{lem}[{\cite[Prop.\,2.2.7]{Lo17b}}]\label{lem_rees_G}
The category $\Rees^G_X$ of Rees $G$-bundles is equivalent to the category $\Fib(\Rep(G), \Rees_X)$ of fiber functors $\Rep(G) \to \Rees_X$.
\end{lem}

For a \emph{filtered $G$-bundle}, it is defined as a pair $(E,\sigma)$, where $E$ is a $G$-bundle and $\sigma: X \to E/P$ is reduction of structure group for some parabolic subgroup $P \subseteq G$. Similar to Lemma \ref{lem_rees_G}, we~have the following equivalent description:

\begin{lem}[{\cite[\S3]{Lo17a}}]\label{lem_fil_G}
The category $\Filt_X^G$ of filtered $G$-bundles is equivalent to the category $\Fib(\Rep(G), \Filt_X)$ of fiber functors $\Rep(G) \to \Filt_X$.
\end{lem}

Combining Lemma \ref{lem_rees_fil}, \ref{lem_rees_G} and \ref{lem_fil_G} altogether
\begin{center}
\begin{tikzcd}
\Rees^G_X \arrow[rr, leftrightarrow, "\text{Lemma \ref{lem_rees_G}}"] \arrow[dd, dotted, leftrightarrow] & & \Fib(\Rep(G), \Rees_X) \arrow[dd, leftrightarrow, "\text{Lemma \ref{lem_rees_fil}}"] \\
& & \\
\Filt^G_X \arrow[rr, leftrightarrow, "\text{Lemma \ref{lem_fil_G}}"] & & \Fib(\Rep(G), \Filt_X),
\end{tikzcd}
\end{center}
we see that the category $\Rees^G_X$ of Rees $G$-bundles is equivalent to the category $\Filt^G_X$ of filtered $G$-bundles. The approach of Tannakian categories not only gives an equivalent description of $G$-bundles, Rees $G$-bundles and filtered $G$-bundles theoretically, but also provides an effective way to work on specific groups.

Now we take $G$ to be orthogonal/symplectic groups, and
denote by $G \hto \GL(W)$ the standard representation. Therefore, an orthogonal/symplectic bundle $E$ is equivalent to a pair $(V, \langle \,, \, \rangle)$, where $V:=E \times_G W$ is the associated bundle and $\langle \,, \, \rangle$ is a nondegenerate symmetric/skew-symmetric bilinear form. Since the base variety $X$ is of arbitrary dimension, it is more flexible to work with orthogonal/symplectic sheaf (see also \cite[\S 5]{GS03}).
\begin{defn}
An \emph{orthogonal/symplectic sheaf} is a pair $(V, \langle \,, \, \rangle)$ such that
\begin{enumerate}
\item $V$ is a torsion free coherent sheaf over $X$,

\item $\langle \,, \, \rangle:V\otimes V\to \sO_X$ is a symmetric/skew-symmetric $\sO_X$-bilinear form and nondegenerate on $U \subseteq X$, where $U$ is an open subset contained in the locally free locus of $V$ with $\codim(X \backslash U)\geq 2$,

\item $\det(V)^2 \cong \mathcal{O}_X$.
\end{enumerate}
\end{defn}
The pairing induces an $\sO_X$-linear morphism $\langle \,, \, \rangle: V\to V^{\vee}$, which is an isomorphism over $U$.
\begin{defn}
Let $(V,\langle \,, \, \rangle)$ be an orthogonal/symplectic sheaf and let $F\subset V$ be a saturated subsheaf.
\begin{enumerate}
\item We call $F^{\perp}:=\ker(V\To{\langle \,, \, \rangle} V^{\vee}\to F^{\vee})$ the \emph{orthogonal complement} of $F$.

\item $F$ is \emph{isotropic} if $\langle \,, \, \rangle|_{F\otimes F}=0$. It is easy to check that $F$ is isotropic if and only if $F\subset F^{\perp}$.
\end{enumerate}
\end{defn}

A filtered orthogonal/symplectic bundle is a pair $(E,\sigma)$, where $\sigma: X \to E/P$ is a reduction of structure group. By Lemma \ref{lem_fil_G}, it is equivalent to a fiber functor $\Rep(G) \to \Filt_X$. Choosing the standard representation $G \hto \GL(W)$, reductions of structure group of $E$ correspond to filtrations of isotropic subbundles of $V:=E \times_G W$ and we refer the reader to \cite[\S 6]{CS21} and \cite[\S 5]{Ya22} for more details. Then a sheaf version is given as follows.

\begin{defn}
Let $(V,\langle \,, \, \rangle)$ be an orthogonal/symplectic sheaf. An \emph{orthogonal/symplectic filtration} of $(V,\langle \,, \, \rangle)$ is a filtration of saturated subsheaves
\begin{align*}
\Filt: 0=\Fil_m \subset \Fil_{m-1} \subset \cdots \subset \Fil_0=V
\end{align*}
such that $\Fil_{i}= \Fil_{m-i}^\perp$. Furthermore, a filtration of saturated subsheaves will be called a \emph{saturated filtration} for simplicity.
\end{defn}

\begin{rem}
Let $(V, \langle \,, \, \rangle)$ be an orthogonal/symplectic sheaf together with an orthogonal/symplectic filtration $\Filt$. There is a natural orthogonal/symplectic structure on the graded bundle
$\Gr_{\Filt}(V):=\mathop{\bigoplus}_{i=0}^{m-1}\sfrac{\Fil_i} {\Fil_{i+1}}$ induced by that on
\begin{align*}
\frac{\Fil_i}{\Fil_{i+1}} \oplus \frac{\Fil_{m-i-1}}{\Fil_{m-i}} \cong \frac{\Fil_{i}}{\Fil_{i+1}} \oplus \frac{\Fil^\perp_{i+1}}{\Fil^\perp_{i}},
\end{align*}
which is denoted by $\overline{\langle \,, \, \rangle}$. Then, we~obtain a graded orthogonal/symplectic sheaf $(\Gr_{\Filt}(V), \overline{\langle \,, \, \rangle})$. Moreover, in the case of orthogonal/symplectic bundles, the graded orthogonal/symplectic bundle $(\Gr_{\Filt}(V), \overline{\langle \,, \, \rangle})$ we obtain is exactly the restriction of the corresponding orthogonal/symplectic Rees bundle to $t=0$ as given in Remark~\ref{rem_graded}.
\end{rem}

\Subsection{Logarithmic Orthogonal/Symplectic $\lambda$-connections}\label{subsect_log_lambda_conn}

\begin{defn}\label{defn_lambda_conn}
Let $(V,\langle \,, \, \rangle)$ be an orthogonal/symplectic sheaf on $X$. Let $\lambda\in k$. A \emph{logarithmic orthogonal/symplectic \lconnection} is a $k$-linear map
\[
\nabla: V \to V\otimes \Omega_X(\log D)
\]
such that on each open subset $U \subseteq X$, it satisfies
\begin{enumerate}
\item \emph{Leibniz rule}:\vspace*{-3pt}
\begin{align*}
\nabla (fa)=\lambda a\otimes df+f\nabla(a),
\end{align*}
where $f\in \mathcal{O}_X(U)$, $a \in V(U)$ and $d: \mathcal{O}_X \to \mathcal{O}_X$ is the exterior differential;

\item \emph{Compatibility}:\vspace*{-3pt}
\begin{align*}
\langle \nabla(a),b \rangle+\langle a,\nabla(b) \rangle=\lambda d\langle a,b \rangle,
\end{align*}
where $a,b \in V(U)$.
\end{enumerate}
In this paper, a logarithmic orthogonal/symplectic \lconnection will be called an \emph{orthogonal/symplectic \lconnection} for simplicity.
\end{defn}

\begin{rem}
Let $(V,\langle \,, \, \rangle,\nabla)$ be an orthogonal/symplectic sheaf together with an orthogonal/symplectic \lconnection. 
\begin{enumerate}
\item When $\lambda=1$, we~say that $\nabla$ is an \emph{orthogonal/symplectic connection}.

\item When $\lambda=0$, $\nabla$ is called an \emph{orthogonal/symplectic Higgs field}. In~this case, we~prefer to use the notation $\theta := \nabla$, and the pair $(V,\theta)$ is called an \emph{orthogonal/symplectic Higgs bundle}. Here we abuse the terminology by ignoring the integrability condition on $\theta$.
\end{enumerate}
\end{rem}

Now we introduce Hodge filtrations for orthogonal/symplectic \lconnections.
\begin{defn}\label{defn_na_Hod_fil}
Given a triple $(V,\langle \,, \, \rangle,\nabla)$ as above, let $\Filt$\vspace*{-3pt}
\begin{align*}
\Filt: 0=\Fil_m \subset \Fil_{m-1} \subset \cdots \subset \Fil_0=V
\end{align*}
be an orthogonal/symplectic filtration of $V$. We~say that $\Filt $ is an \emph{orthogonal/sym\-plec\-tic Hodge filtration} if $\Filt $ satisfies the \emph{Griffiths transversality condition}:\vspace*{-3pt}
\[
\nabla(\Fil_{i+1})\subset \Fil_{i}\otimes \Omega_X(\log D), \quad i \geq 1.
\]
Such a tuple $(V, \langle \,, \, \rangle, \nabla, \Filt)$ is called a \emph{Hodge filtered orthogonal/symplectic \lconnection}.
\end{defn}

Let $(V, \langle \,, \, \rangle,\nabla, \Filt)$ be a Hodge filtered orthogonal/symplectic \lconnection. The Rees construction gives a sheaf\vspace*{-3pt}
\begin{align*}
\mathcal{V}= \sum t^{-i} \Fil_{i} \otimes \mathcal{O}_{X \times \mathbb{A}^1}
\end{align*}
on $X \times \mathbb{A}^1$, and $t \nabla$ is well-defined on $\mathcal{V}$. Taking the limit $\lim_{t \to 0} (V,t\nabla)$, we~obtain a well-defined Higgs field $\theta$ on $\Gr_{\Filt}(V)$. In~conclusion, we~obtain an orthogonal/symplectic Higgs sheaf $(\Gr_{\Filt}(V), \overline{\langle \,, \rangle}, \theta)$ such that
\begin{align*}
\theta( \sfrac{\Fil_i}{\Fil_{i+1}} ) \subseteq \Psfrac{\Fil_{i-1}}{\Fil_{i}} \otimes \Omega_X(\log D), \ i \geq 1.
\end{align*}
Note that $\Gr_{\Filt}(V)$ is torsion-free because the filtration we choose is saturated. This observation gives the following definition.

\begin{defn}
A \emph{graded orthogonal/symplectic Higgs sheaf} is an orthogonal/symplectic Higgs sheaf $(E, \langle \,, \, \rangle,\theta)$ together with a decomposition $E=\mathop{\bigoplus}_{i=1}^{m} E_i$ such that
\begin{enumerate}
\item $\theta(E_i)\subset E_{i+1}\otimes \Omega_X(\log D)$ for $ 1\leq i\leq m-1$, and $\theta(E_m)=0$,

\item $\langle \,, \, \rangle|_{E_i\otimes E_j}$ is nondegenerate for $i+j=m+1$ over a nonempty open subset in the locally free locus whose complement is of codimension $\geq 2$,

\item $\langle \,, \, \rangle|_{E_i\otimes E_j}=0$ for $i+j\neq m+1$.
\end{enumerate}
\end{defn}

Summing up, as in the case of vector bundles \cite{S10}, one gets a graded orthogonal/symplectic Higgs sheaf from a Hodge filtered orthogonal/symplectic \lconnection:

\begin{lem}
Let $(V, \langle \,, \, \rangle,\nabla, \Filt)$ be a Hodge filtered orthogonal/symplectic \lconnection. It induces a graded orthogonal/symplectic Higgs sheaf $(E,\overline{\langle \,, \, \rangle},\theta)$, where $E:=\Gr_{\Filt}(V)$.
\end{lem}

We would like to remind the reader that in the above lemma, the \lconnection $\nabla$ could be a connection or a Higgs field.

\subsection{Stability condition and Harder--Narasimhan filtration}\label{subsect_stab_cond}
We fix an ample line bundle $L$ over $X$. Let $V$ be a torsion free sheaf on $X$. Define
\[
\mu_L(V):=\deg_L(V)/\rk(V)
\]
to be the slope of $V$ with respect to $L$. In~this paper, if there is no ambiguity, we~omit the subscript $L$ and use the notation $\mu$ and $\deg$ for slope and degree respectively. In~this subsection, a triple $(V,\langle \,, \, \rangle,\nabla)$ is an orthogonal/symplectic sheaf $(V,\langle \,, \, \rangle)$ together with an orthogonal/symplectic \lconnection~$\nabla$ unless stated otherwise.

\begin{defn}\label{defn_stab}
A triple $(V,\langle \,, \, \rangle,\nabla)$ is \emph{semistable} (\resp \emph{stable}) if for any nontrivial $\nabla$-invariant isotropic subsheaf $W \subset V$, we~have $\mu(W)\leqslant 0$ (\resp $<$).
\end{defn}

\begin{rem}
In this remark, we~only focus on the case of curves. In~characteristic zero, the slope stability condition of orthogonal/symplectic bundles $(V, \langle \,, \, \rangle)$ is equivalent to the Ramanathan's stability condition on the corresponding principal bundles~$\mathcal{V}$ \cite{Ra75,Ra96a,Ra96b}. We~briefly review the correspondence as follows, and the Ramanathan stability condition will be called $R$-stability condition for simplicity, where ``$R$'' stands for Ramanathan.

A $G$-bundle $\mathcal{V}$ is \emph{$R$-semistable} (\resp \emph{$R$-stable}) if for any proper parabolic subgroup $P \subseteq G$, any reduction of structure group $\sigma: X \to \mathcal{V}/P$ and any dominant character $\chi: P \to \mathbb{G}_m$, we~have $\deg \chi_* \mathcal{V}_\sigma \leq 0$ (\resp $<$), where $\mathcal{V}_\sigma := \mathcal{V} \times_{\sigma} X$ is a $P$-bundle on $X$ and $\chi_* (\mathcal{V}_\sigma)$ is a line bundle on $X$. When we consider the case of orthogonal/symplectic groups, reductions of structure group correspond exactly to filtrations of isotropic subbundles. Following this idea, the equivalence of the stability conditions for symplectic groups is proved in \cite[App.]{KSZ21}, and the same argument holds for orthogonal groups.

In positive characteristic, the $R$-stability condition of $\mathcal{V}$ is equivalent to the stability condition of the associated bundle under a low weight representation $G \hto \SL(W)$ \cite[\S 2]{BP03}. In~our case, the representation we take is the standard one, which is obvious of low weight. Therefore, the stability condition of orthogonal/symplectic bundles considered in this paper is equivalent to the Ramanathan's stability condition for principal bundles in positive characteristic.

Furthermore, the above result also holds for \lconnections. More precisely, let $(V, \langle \,, \, \rangle, \nabla)$ be an orthogonal/symplectic \lconnection and denote by $(\mathcal{V},\nabla')$ the corresponding $G$-connection, where $G$ is the orthogonal/symplectic group. Then the slope stability condition of $(V, \langle \,, \, \rangle, \nabla)$ given in Definition \ref{defn_stab} is equivalent to the $R$-stability condition of $(\mathcal{V},\nabla')$. The case of $G=\GL_n$ and $\lambda = 0$ is proved in \cite[Prop.\,5.5]{KSZ23} and the orthogonal/symplectic case with arbitrary $\lambda$ can be proved in the same way.
\end{rem}

In the following, we~prove that the Harder--Narasimhan filtration (for slope stability condition) of an orthogonal/symplectic connection is exactly the Harder--Narasimhan filtration for the corresponding connection by forgetting the orthogonal/symplectic structure (Proposition \ref{prop_HN_fil}). As a direct result, we~obtain the maximal destabilizer of $(V, \langle \,,\, \rangle, \nabla)$ is isotropic (Corollary \ref{cor_max_destab_iso}), which will be used frequently later on.

\begin{prop}\label{prop_HN_fil}
Given a triple $(V,\langle \,, \, \rangle,\nabla)$, let
\[
\HN:0=V_0\subset V_1 \subset \cdots \subset V_t=V
\]
be the Harder--Narasimhan filtration of $(V, \nabla)$, that is, we~forget the orthogonal/symplectic structure of $V$. Then $\HN$ is an orthogonal/symplectic filtration.
\end{prop}

\begin{proof}
In fact,
\[
\HN^\perp:0=V_t^\perp\subset V_{t-1}^\perp \subset \cdots \subset V_0^\perp=V
\]
is again a Harder--Narasimhan filtration of $(V, \nabla)$. By the uniqueness of the Harder--Narasimhan filtration, we~have $V_{t-i}=V_{i}^\perp$.
\end{proof}

As a direct consequence, we~have the following corollary.
\begin{cor}\label{cor_max_destab_iso}
Suppose that $(V,\langle \,, \, \rangle,\nabla)$ is unstable. Then the maximal destabilizer of $(V, \nabla)$ is isotropic. Furthermore, $(V,\langle \,, \, \rangle,\nabla)$ is semistable if and only if $(V, \nabla)$ is semistable.
\end{cor}

\begin{notn}
Let
\begin{align*}
\HN:0=V_0\subset V_1 \subset \cdots \subset V_t=V
\end{align*}
be the Harder--Narasimhan filtration of $(V,\langle \,, \, \rangle,\nabla)$. Define $\mu_{\max}(V)$ (\resp $\mu_{\min}(V)$) to be the slope of $V_1 / V_0$ (\resp $V_t/V_{t-1}$). We~also refer the reader to \cite[Def.\,1.3.2]{HL97} for a related notation. Moreover, the orthogonal/symplectic \lconnection $(V,\langle \,, \, \rangle,\nabla)$ can be replaced by torsion free sheaves, orthogonal/symplectic sheaves, \lconnections and etc.
\end{notn}

\skpt
\begin{defn} \label{gr-semistable}
Let $\Filt$ be an orthogonal/symplectic Hodge filtration on $(V, \langle \,,\, \rangle,\nabla)$. We~say that $\Filt$ is \emph{\grsemistable} if the associated graded orthogonal/sym\-plec\-tic Higgs sheaf $(\Gr_{\Filt}(V), \overline{\langle \,, \, \rangle},\theta)$ is semistable. Moreover, $(V, \langle \,, \, \rangle, \nabla)$ is \emph{\grsemistable} if there exists a \grsemistable filtration $\Filt$ of $(V, \langle \,,\, \rangle,\nabla)$.
\end{defn}

One may ask that whether there always exists a \grsemistable filtration on a semistable triple $(V, \langle \,,\, \rangle,\nabla)$, where $\nabla$ is a connection? We will give an answer to this problem in the next sections.

\section{Quasi \texorpdfstring{\protect\grsemistable}{grsemistable} filtration: odd rank}\label{sect_odd}

Let $(V, \langle \,, \, \rangle, \nabla)$ be an orthogonal sheaf $(V, \langle \,, \, \rangle)$ together with an orthogonal \lconnection $\nabla$, where $V$ is a torsion free sheaf of odd rank. In~this section, we~introduce a special type of orthogonal Hodge filtrations, which is called \emph{quasi \grsemistable filtrations} (Definition \ref{defn_gr_semi_odd}), and prove that the existence of \grsemistable filtrations is equivalent to the existence of quasi \grsemistable filtrations of $(V, \langle \,, \, \rangle, \nabla)$ (Theorem~\ref{thm_gr_and_quasi_gr_odd}).

\Subsection{Definition and Basic Properties}

\begin{defn}[{Quasi \grsemistable filtration in odd rank case}]\label{defn_gr_semi_odd}
Let
\[
\Filt: 0= L_0 \subsetneq L_1 \subsetneq \cdots \subsetneq L_m \subsetneq L_m^\perp\subsetneq\cdots \subsetneq L_1^\perp\subsetneq L_0^\perp= V
\]
be an orthogonal Hodge filtration of $(V, \langle \,, \, \rangle, \nabla)$. We~say that $\Filt$ is \emph{quasi \grsemistable} if it satisfies the following two conditions:

\begin{enumerate}\renewcommand{\theenumi}{\Roman{enumi}}
\item\label{cond:I}
For any saturated Griffiths transverse filtration
\[
0=W_0\subset W_1 \subset \cdots\subset W_m
\]
such that $L_{i-1}\subset W_i\subset L_i$, we~have
\[
\sum_{0\leq i\leq m}\deg(W_i)\leqslant\sum_{0\leq i\leq m}\deg(L_i).
\]

\item\label{cond:II}
For any saturated Griffiths transverse filtration
\[
0\subset S_1 \subset \cdots\subset S_m\subset S_{m+1}\subset T_m^\perp\subset\cdots \subset T_1^\perp\subset V
\]
such that
\[
\deg(S_{m+1})>\sum_{0\leq i\leq m}\big(\deg(L_i)-\deg(S_i)\big)+\sum_{0\leq i\leq m}\big(\deg(L_i)-\deg(T_i)\big),
\]
we have
\[
\nabla(S_{m+1})\subseteq S_{m+1}^\perp \otimes\Omega_X(\log D),
\]
where $L_{i-1}\subset S_i \subset T_i \subset L_i,\ L_m\subset S_{m+1}\subset L_m^\perp$ and $S_{m+1}$ is isotropic.
\end{enumerate}
Moreover, we~say that $(V,\langle \,, \, \rangle,\nabla)$ is \emph{quasi \grsemistable} (or $\nabla$ is a \emph{quasi \grsemistable connection}) if there exists a quasi \grsemistable filtration on $V$.
\end{defn}

\begin{examp}\label{smallm}
In this example, we~briefly explain the quasi gr-semistability condition for small $m$.

\begin{itemize}
\item $m=0$.
The trivial filtration
\[
\Filt:0\subset V
\]
is quasi \grsemistable if and only if $\nabla(W)\subset W^\perp\otimes\Omega_X(\log D)$ holds for any isotropic saturated subsheaf $W$ of $\deg(W)>0$.

\item $m=1$.
An orthogonal Hodge filtration $0\subset N\subset N^\perp\subset V$ is quasi \grsemistable if and only if it satisfies the following two conditions:
\begin{enumerate}
\item $\mu_{\min}(N)\geq 0$;
\item for any saturated Griffiths transverse filtration $0\!\subset\! S_1\!\subset\! S_2\!\subset\! T_1^\perp$ such that
\begin{itemize}
\item[$\scriptscriptstyle\bullet$] $S_1\subset T_1\subset N,\ N\subset S_2\subset N^\perp $,
\item[$\scriptscriptstyle\bullet$] $S_2$ is isotropic,
\item[$\scriptscriptstyle\bullet$] $\deg(S_2)>\deg(N)-\deg(S_1)+\deg(N)-\deg(T_1)$,
\end{itemize}
we have $\nabla(S_2)\subset S_2^\perp\otimes \Omega_X(\log D)$.
\end{enumerate}
\end{itemize}
When $m=0$, quasi gr-semistability is easy to check. However, when $m \geq 1$, the stability condition becomes much more complicated. The above discussion will be used in Section~\ref{sect_small_rank}.
\end{examp}

\begin{rem}\label{explainregular}
In this remark, we~briefly explain Condition \eqref{cond:I} and \eqref{cond:II} in terms of graded Higgs sheaves. Given $(V, \langle \,, \, \rangle, \nabla, \Filt)$, let $(\Gr_{\Filt}, \overline{\langle \,, \, \rangle}, \theta)$ be the corresponding graded orthogonal Higgs sheaf, where
\begin{align*}
\Gr_{\Filt}(V)=\mathop{\bigoplus}_{i=1}^{m} \frac{L_i}{L_{i-1}}\oplus \frac{L_m^\perp}{L_m}\oplus \mathop{\bigoplus}_{i=1}^{m} \frac{L_{m-i}^\perp}{L_{m-i+1}^\perp}.
\end{align*}

\begin{enumerate}
\item The filtration
\[
0=W_0\subset W_1 \subset \cdots\subset W_m
\]
in Condition \eqref{cond:I} gives a Higgs subsheaf of $\Gr_{\Filt}(V)$
\[
\alpha(W_\sbullet):=\mathop{\bigoplus}_{i=1}^{m} \frac{W_i}{L_{i-1}}\oplus \frac{L_m^\perp}{L_m}\oplus \mathop{\bigoplus}_{i=1}^{m} \frac{L_{m-i}^\perp}{L_{m-i+1}^\perp}.
\]
Then the inequality
\[
\sum_{0\leq i\leq m}\deg(W_i)\leqslant\sum_{0\leq i\leq m}\deg(L_i).
\]
is equivalent to $\deg\big(\alpha(W_\sbullet)\big)\leq 0$. 
\item The filtration
\[
0\subset S_1 \subset \cdots\subset S_m\subset S_{m+1}\subset T_m^\perp\subset\cdots \subset T_1^\perp\subset V
\]
in Condition \eqref{cond:II} defines an isotropic Higgs subsheaf of $\Gr_{\Filt}(V)$
\[
\beta(S_\sbullet,T_\sbullet):=\mathop{\bigoplus}_{i=1}^{m} \frac{S_i}{L_{i-1}}\oplus \frac{S_{m+1}}{L_m}\oplus \mathop{\bigoplus}_{i=1}^{m} \frac{T_{m-i+1}^\perp}{L_{m-i+1}^\perp}.
\]
Then the inequality
\[
\deg(S_{m+1})>\sum_{0\leq i\leq m}\big(\deg(L_i)-\deg(S_i)\big)+\sum_{0\leq i\leq m}\big(\deg(L_i)-\deg(T_i)\big)
\]
is equivalent to $\deg\big(\beta(S_\sbullet,T_\sbullet)\big)>0$.
\end{enumerate}

Summing up, $\Filt$ is quasi \grsemistable if and only if the following two conditions hold:
\begin{enumerate}
\item $\deg\big(\alpha(W_\sbullet)\big)\leq 0$;

\item if $\deg\big(\beta(S_\sbullet,T_\sbullet)\big)>0$, then we have $\nabla(S_{m+1})\subseteq S_{m+1}^\perp \otimes\Omega_X(\log D)$.
\end{enumerate}
\end{rem}

\begin{prop}\label{ssregular}
If an orthogonal Hodge filtration $\Filt$ of $(V,\langle \,, \, \rangle,\nabla)$ is \grsemistable, then $\Filt$ is quasi \grsemistable.
\end{prop}

\begin{proof}
The semistability of $\big(\Gr_{\Filt}(V), \overline{\langle \,, \, \rangle}, \theta\big)$ implies that $\deg(W)\leq 0$ holds for any isotropic Higgs subsheaf of $\big(\Gr_{\Filt}(V),\theta\big)$. By Prop \ref{prop_HN_fil}, $\deg(W)\leq 0$ holds for any Higgs subsheaf of $\big(\Gr_{\Filt}(V),\theta\big)$. Hence for any saturated Griffiths transverse filtration given in Condition \eqref{cond:I} and \eqref{cond:II}, we~have $\deg\big(\alpha(W_\sbullet)\big)\leq 0$ and $\deg\big(\beta(S_\sbullet,T_\sbullet)\big)\leq0$. Then the first condition in Remark \ref{explainregular} is satisfied automatically, and the second condition also holds because there is no $\beta(S_\sbullet,T_\sbullet)$ such that $\deg( \beta(S_\sbullet,T_\sbullet)) > 0$.
\end{proof}

Remark \ref{explainregular} and Proposition \ref{ssregular} is actually the starting point how we introduce the definition of quasi gr-semistability. The idea is simple and we briefly state as follows. Given an orthogonal/symplectic Hodge filtration $\Filt$ on a semistable triple $(V, \langle \,,\, \rangle, \nabla)$, we~hope that the resulting graded orthogonal/symplectic Higgs sheaf $\big(\Gr_{\Filt}(V), \overline{\langle \,, \, \rangle}, \theta\big)$ is semistable. More precisely, the slope of any $\theta$-invariant isotropic subsheaf of $\Gr_{\Filt}(V)$ should be smaller than or equal to zero. Then we interpret the above inequalities in terms of specific filtrations on $V$ and find a necessary condition, which implies the gr-semistability condition. This condition is called quasi gr-semistability condition and is given in Definition \ref{defn_gr_semi_odd} and explained in Remark \ref{explainregular}. Therefore, the proof of Proposition \ref{ssregular} follows automatically. On the other hand, the most surprising fact is that the condition we introduce (Definition \ref{defn_gr_semi_odd}) is not only a necessary condition, but also a sufficient one, which will be proved in the next subsection.

\begin{prop}\label{positive}
Let
\[
\Filt: 0= L_0 \subsetneq L_1 \subsetneq \cdots \subsetneq L_m \subsetneq L_m^\perp\subsetneq\cdots \subsetneq L_1^\perp\subsetneq L_0^\perp= V
\]
be a nontrivial quasi \grsemistable filtration of $(V, \langle \,, \, \rangle, \nabla)$. Then we have $\deg(L_i)\geq 0$. Moreover, $\deg(F)\leq \deg(L_1)$ holds for any subsheaf $F\subset L_1$.
\end{prop}

\begin{proof}
For a fixed $1\leq i_0\leq m$, consider the following Griffiths transverse filtration
\[
0=W_0\subset W_1\subset \cdots\subset W_m,
\]
where $W_i= L_{i-1}$ for $1\leq i\leq i_0$ and $W_i= L_i$ for $i_0+1\leq i\leq m$. Then the inequality
\[
\sum_{0\leq i\leq m}\deg(W_i)\leqslant\sum_{0\leq i\leq m}\deg(L_i)
\]
implies that $\deg(L_{i_0})\geq 0$. Similarly, let $W_1=F$ and $W_i=L_i$ for $2\leq i\leq m$, and we get $\deg(F)\leq \deg(L_1)$.
\end{proof} 

\subsection{The existence of \texorpdfstring{\protect\grsemistable}{grsemistable} filtration}\label{subsect_exist}

In this subsection, we~will show that given a semistable triple $(V, \langle \,, \, \rangle, \nabla)$, the quasi gr-semistability of $(V, \langle \,, \, \rangle, \nabla)$ implies the existence of \grsemistable filtrations (Theorem \ref{main thm}). Combining with Proposition \ref{ssregular}, we~obtain the main result (Theorem \ref{thm_gr_and_quasi_gr_odd}) in this section. First, let us begin with a lemma about the maximal destabilizer of a graded orthogonal Higgs sheaf.

\begin{lem}[{\cite[Lem.\,A.6]{LSZ19}}]\label{destabilizer}
Let $(E=\mathop{\bigoplus}_{i=1}^{m} E_i,\ \langle \,, \, \rangle,\ \theta)$ be an unstable graded orthogonal Higgs sheaf. Let $M\subset E$ be the maximal destabilizer of $(E,\theta)$. Then $M$ is an isotropic saturated graded Higgs subsheaf, that is, $M=\mathop{\bigoplus}_{i=1}^{m} M_i$ with $M_i=M\cap E_i$ and the Higgs field is induced by $\theta$.
\end{lem}

Now let $\Filt$
\[
\Filt: 0= L_0 \subsetneq L_1 \subsetneq \cdots \subsetneq L_m \subsetneq L_m^\perp\subsetneq\cdots \subsetneq L_1^\perp\subsetneq L_0^\perp= V
\]
be a quasi \grsemistable filtration of $(V,\langle \,, \, \rangle,\nabla)$, and we consider the associated orthogonal graded Higgs sheaf $(\Gr_{\Filt}(V), \overline{\langle \,, \, \rangle},\theta)$, where
\[
\Gr_{\Filt}(V)=\mathop{\bigoplus}_{i=1}^{m} \frac{L_i}{L_{i-1}}\oplus \frac{L_m^\perp}{L_m}\oplus \mathop{\bigoplus}_{i=1}^{m} \frac{L_{m-i}^\perp}{L_{m-i+1}^\perp}.
\]
If $(\Gr_{\Filt}(V),\theta)$ is not semistable, then we have a maximal destabilizer and denote it by $M_{\Filt}$. By Lemma \ref{destabilizer} and Corollary \ref{cor_max_destab_iso}, we~have an isotropic subsheaf
\[
M_{\Filt}=\mathop{\bigoplus}_{i=1}^{m} \frac{M_i}{L_{i-1}}\oplus \frac{M_{m+1}}{L_m}\oplus \mathop{\bigoplus}_{i=1}^{m} \frac{N_{m-i+1}^\perp}{L_{m-i+1}^\perp}.
\]
Moreover, note that the induced bilinear form on $M_i \otimes N_i^{\perp}$ is trivial, and thus $L_{i-1} \subseteq M_i \subseteq N_i \subseteq L_i$. Moreover, it is easy to see that $M_{m+1}$ is an isotropic subsheaf. We~construct a new filtration $\xi(\Filt)$ from $\Filt$ as follows:
\[
\xi(\Filt):0=M_0\subset M_1\subset\cdots\subset M_{m+1}\subset M_{m+1}^\perp\subset\cdots\subset M_1^\perp\subset M_0^\perp=V.
\]
Moreover, $\xi$ can be regarded as an operator on the set of Griffiths transverse filtrations.

\begin{lem}\label{preserve}
Let $\Filt$ be a quasi \grsemistable filtration of $(V, \langle \,, \, \rangle, \nabla)$. Then $\xi(\Filt)$ is a (orthogonal) Hodge filtration. Moreover, $\xi(\Filt)$ is also quasi \grsemistable.
\end{lem}

\begin{proof}
First, we~will show that $\xi(\Filt)$ is a Hodge filtration, that is, $\xi(\Filt)$ satisfies the Griffiths transversality condition. It is enough to show that
\[
\nabla(M_{m+1})\subset M_{m+1}^\perp\otimes\Omega_X(\log D).
\]
This is true by Remark \ref{explainregular} since $\Filt$ is quasi \grsemistable and $\deg(M_{\Filt})>0$.

Next, we~will check the quasi gr-semistability of $\xi(\Filt)$.

\eqref{cond:I}
Let
\[
0=W_0\subset W_1 \subset \cdots\subset W_{m+1}
\]
be a given saturated Griffiths transverse filtration such that $M_{i-1}\subset W_i\subset M_i$. We~need to show that
\[
\sum_{0\leq i\leq m+1}\deg(W_i)\leqslant\sum_{0\leq i\leq m+1}\deg(M_i).
\]
Consider the following Higgs subsheaf $\sigma(W_\sbullet)$ of $\Gr_{\Filt}(V)$:
\[
\sigma(W_\sbullet):=\mathop{\bigoplus}_{i=1}^{m} \frac{W_i+L_{i-1}}{L_{i-1}}\oplus\frac{W_{m+1}+L_m}{L_m} \oplus \mathop{\bigoplus}_{i=1}^{m} \frac{N_{m-i+1}^\perp}{L_{m-i+1}^\perp}.
\]
Clearly, $\sigma(W_\sbullet)$ is a Higgs subsheaf of $M_{\Filt}$. Since $M_{\Filt}$ is the maximal destabilizer, we~have
\begin{equation}
\deg(M_{\Filt})-\deg(\sigma(W_\sbullet))\geq 0
\end{equation}
On the other hand,
\begin{align*}
\deg&(M_{\Filt})-\deg(\sigma(W_\sbullet))=\sum_{0\leq i\leq m+1}\deg(M_i)-\sum_{0\leq i\leq m+1}\deg(W_i+L_{i-1})\\
&=\sum_{0\leq i\leq m+1}\deg(M_i)-\sum_{0\leq i\leq m+1}\big(\deg(W_i)+\deg(L_{i-1})-\deg(W_i\cap L_{i-1})\big)\\
&=\sum_{0\leq i\leq m+1}\big(\deg(M_i)-\deg(W_i)\big)-\sum_{0\leq i\leq m+1}\big(\deg(L_{i-1})-\deg(W_i\cap L_{i-1})\big).
\end{align*}
Therefore,
\begin{equation}\label{increasing}
\sum_{0\leq i\leq m+1}\big(\deg(M_i)-\deg(W_i)\big)\geq \sum_{0\leq i\leq m+1}\big(\deg(L_{i-1})-\deg(W_i\cap L_{i-1})\big)\geq 0.
\end{equation}
The last inequality holds because $\Filt$ is quasi \grsemistable. Condition \eqref{cond:I} in Definition~\ref{defn_gr_semi_odd} is satisfied.

\eqref{cond:II}
Let
\[
\Filt_{ST}:0\subset S_1 \subset \cdots\subset S_{m+1}\subset S_{m+2}\subset T_{m+1}^\perp\subset\cdots \subset T_1^\perp\subset V
\]
be a given saturated Griffiths transverse filtration such that
\begin{itemize}
\item $M_{i-1}\subset S_i \subset T_i \subset M_i,\ M_{m+1}\subset S_{m+2}\subset M_{m+1}^\perp$,
\item $S_{m+2}$ is isotropic,
\item the inequality
\begin{equation}\label{inequality}
\deg(S_{m+2})>\sum_{0\leq i\leq m+1}\big(\deg(M_i)-\deg(S_i)\big)+\sum_{0\leq i\leq m+1}\big(\deg(M_i)-\deg(T_i)\big)
\end{equation}
is satisfied.
\end{itemize}
We need to show that $\nabla(S_{m+2})\subset S_{m+2}^\perp\otimes \Omega_X(\log D)$. Consider the following filtration:
\[
\Filt_{\tilde{S}\tilde{T}}:=0\subset \tilde{S}_1\subset \cdots\subset \tilde{S}_m\subset \tilde{S}_{m+1}\subset \tilde{T}_m^\perp\subset\cdots \subset \tilde{T}_1^\perp\subset V,
\]
where
\[
\tilde{S}_i=S_{i+1}\cap L_i,\quad \tilde{T}_i=T_{i+1}\cap L_i,\quad \text{for $1\leq i\leq m$},\quad\text{and}\quad\tilde{S}_{m+1}=S_{m+2}.
\]
Then $L_{i-1}\subset \tilde{S}_i\subset \tilde{T}_i\subset L_i$ for $1\leq i\leq m$, $L_m\subset\tilde{S}_{m+1}\subset L_m^\perp$ and $\tilde{S}_{m+1}$ is isotropic. Moreover, we~have
\[
\deg(\tilde{S}_{m+1})>\sum_{0\leq i\leq m}\big(\deg(L_i)-\deg(\tilde{S}_i)\big)+\sum_{0\leq i\leq m}\big(\deg(L_i)-\deg(\tilde{T}_i)\big).
\]
In fact, by the inequality \ref{increasing} and \ref{inequality}, we~have
\begin{equation}
\begin{split}
\deg(\tilde{S}_{m+1})&=\deg(S_{m+2})\\
&>\sum_{0\leq i\leq m+1}\big(\deg(M_i)-\deg(S_i)\big)+\sum_{0\leq i\leq m+1}\big(\deg(M_i)-\deg(T_i)\big)\\
&\geq \sum_{0\leq i\leq m+1}\big(\deg(L_{i-1})-\deg(S_i\cap L_{i-1})\big) \\
& \hspace*{2.7cm}+\sum_{0\leq i\leq m+1}\big(\deg(L_{i-1})-\deg(T_i\cap L_{i-1})\big)\\
&=\sum_{0\leq i\leq m}\big(\deg(L_i)-\deg(\tilde{S}_i)\big)+\sum_{0\leq i\leq m}\big(\deg(L_i)-\deg(\tilde{T}_i)\big).
\end{split}
\end{equation}
Then, by the quasi gr-semistability of $\Filt$, we~have
\[
\nabla(\tilde{S}_{m+1})\subset\tilde{S}_{m+1}^\perp\otimes\Omega_X(\log D),
\]
which implies
\[
\nabla(S_{m+2})\subset S_{m+2}^\perp\otimes \Omega_X(\log D).\qedhere
\]
\end{proof}

Now we will state and prove the main result in this section.
\begin{thm}\label{main thm}
Let $(V,\langle \,, \, \rangle,\nabla)$ be semistable and quasi \grsemistable with odd rank. There exists a \grsemistable filtration on $(V,\langle \,, \, \rangle,\nabla)$.
\end{thm}

We need the following two lemmas to prove this theorem.
\begin{lem}\label{decrease}
Let $\Filt$ be a quasi \grsemistable filtration of $(V, \langle \,, \, \rangle, \nabla)$. Assume that the graded orthogonal sheaves $(\Gr_{\Filt}(V),\theta_1)$ and $(\Gr_{\xi(\Filt)}(V),\theta_2)$ are not semistable. Let $M_{\Filt}$ (\resp $M_{\xi(\Filt)}$) be the maximal destabilizer of $(\Gr_{\Filt}(V),\theta_1)$ (resp.\ $\Gr_{\xi(\Filt)}(V),\theta_2)$). Then we have a natural morphism of graded Higgs sheaves
\[
\Phi: M_{\xi(\Filt)}\to M_{\Filt}
\]
such that
\begin{enumerate}
\item $\mu(M_{\xi(\Filt)})\leq \mu(M_{\Filt})$;
\item if $\mu(M_{\xi(\Filt)})= \mu(M_{\Filt})$, the morphism $\Phi$ is injective, then $\rk(M_{\xi(\Filt)})\leq\rk(M_{\Filt})$;
\item if $\big(\mu(M_{\xi(\Filt)},\ \rk(M_{\xi(\Filt)})\big)= \big(\mu(M_{\Filt}),\rk(M_{\Filt})\big)$, then $\Phi$ is an isomorphism on an open subset $U \subset X$ with $\codim(X/U) \geq 2$.
\end{enumerate}
\end{lem}

\begin{proof}
We write $\Filt$ as
\[
\Filt: 0= L_0 \subsetneq L_1 \subsetneq \cdots \subsetneq L_m \subsetneq L_m^\perp\subsetneq\cdots \subsetneq L_1^\perp\subsetneq L_0^\perp= V.
\]
Then,
\[
\Gr_{\Filt}(V)=\mathop{\bigoplus}_{i=1}^{m} \frac{L_i}{L_{i-1}}\oplus \frac{L_m^\perp}{L_m}\oplus \mathop{\bigoplus}_{i=1}^{m} \frac{L_{m-i}^\perp}{L_{m-i+1}^\perp}.
\]
Then, let $M_{\Filt}$ be as follows
\[
M_{\Filt}=\mathop{\bigoplus}_{i=1}^{m} \frac{M_i}{L_{i-1}}\oplus \frac{M_{m+1}}{L_m}\oplus \mathop{\bigoplus}_{i=1}^{m} \frac{N_{m-i+1}^\perp}{L_{m-i+1}^\perp}.
\]
Recall that
\begin{itemize}
\item $L_{i-1} \subset M_i \subset N_i \subset L_i$ for $1 \leq i \leq m$,
\item $M_{m+1}$ is isotropic.
\end{itemize}
Then, $\xi(\Filt)$ is given by
\[
\xi(\Filt):0=M_0\subset M_1\subset\cdots\subset M_{m+1}\subset M_{m+1}^\perp\subset\cdots\subset M_1^\perp\subset M_0^\perp=V,
\]
and similarly,
\[
\Gr_{\xi(\Filt)}(V)=\mathop{\bigoplus}_{i=1}^{m+1} \frac{M_i}{M_{i-1}}\oplus \frac{M_{m+1}^\perp}{M_{m+1}}\oplus \mathop{\bigoplus}_{i=1}^{m+1} \frac{M_{m+1-i}^\perp}{M_{m+2-i}^\perp}.
\]
Therefore, $M_{\xi(\Filt)}$ is given as follows:
\[
M_{\xi(\Filt)}=\mathop{\bigoplus}_{i=1}^{m+1} \frac{X_i}{M_{i-1}}\oplus \frac{X_{m+2}}{M_{m+1}}\oplus \mathop{\bigoplus}_{i=1}^{m+1} \frac{Y_{m+2-i}^\perp}{M_{m+2-i}^\perp},
\]
where
\begin{itemize}
\item $M_{i-1} \subset X_i \subset Y_i \subset M_i$ for $1 \leq i \leq m+1$,
\item $X_{m+2}$ is isotropic.
\end{itemize}

Now we define the following natural morphisms of sheaves:
\begin{equation}
\begin{split}
&\Phi_i:\frac{X_i}{M_{i-1}}\to \frac{M_i}{L_{i-1}},\ 1\leq i\leq m+1,\\
&\Phi_{m+2}:\frac{X_{m+2}}{M_{m+1}}\to 0,\\&\tilde{\Phi}_j:\frac{Y_{m+2-j}^\perp}{M_{m+2-j}^\perp}\to 0,\ 1\leq j\leq m+1,\\
\end{split}
\end{equation}
and we obtain a morphism of graded Higgs sheaves
\[
\Phi: M_{\xi(\Filt)}\to M_{\Filt}.
\]
\begin{enumerate}
\item If $\Phi\neq 0$, we~have
\[
\mu(M_{\xi(\Filt)})\leq \mu(M_{\Filt}).
\]

\item If $\Phi= 0$, then $X_i\subset L_{i-1},\ 1\leq i\leq m+1$. Also, we~have $X_{m+2}\subset L_{m}^\perp$ and $\ Y_{m+2-j}^\perp \subset L_{m-j}^\perp$ for $1\leq j\leq m+1$. In~this case, we~consider the quotient Higgs sheaf
\[
\frac{\Gr_{\Filt}(V)}{M_{\Filt}}=\mathop{\bigoplus}_{i=1}^{m} \frac{L_i}{M_i}\oplus \frac{L_m^\perp}{M_{m+1}}\oplus \mathop{\bigoplus}_{i=1}^{m} \frac{L_{m-i}^\perp}{N_{m-i+1}^\perp},
\]
and define the following natural morphisms:
\begin{equation}
\begin{split}
&\Psi_i:\frac{X_i}{M_{i-1}}\to \frac{L_{i-1}}{M_{i-1}},\qquad 2\leq i\leq m+1,\\
&\Psi_{m+2}:\frac{X_{m+2}}{M_{m+1}}\to \frac{L_m^\perp}{M_{m+1}},\\
&\tilde{\Psi}_j:\frac{Y_{m+2-j}^\perp}{M_{m+2-j}^\perp}\to \frac{L_{m-j}^\perp}{N_{m-j+1}^\perp},\qquad 1\leq j\leq m,\\
&\tilde{\Psi}_{m+1}:\frac{Y_1^\perp}{M_1^\perp}\to 0.
\end{split}
\end{equation}
This construction gives a morphism of graded Higgs sheaves
\[
\Psi: M_{\xi(\Filt)}\to \frac{\Gr_{\Filt}(V)}{M_{\Filt}}.
\]
We \emph{claim} that $\Psi\neq 0$. Then we have
\[
\mu(M_{\xi(\Filt)})\leq \mu(\mathrm{Im}(\Psi)) <\mu(M_{\Filt}).
\]
Note that in this case we get a strict inequality
\[
\mu(M_{\xi(\Filt)})<\mu(M_{\Filt}).
\]
The first statement that $\mu(M_{\xi(\Filt)})\leq \mu(M_{\Filt})$ is proved.

For the \emph{claim}, we~suppose that $\Psi=0$. Then $\sfrac{X_i}{M_{i-1}}=0$ and thus $X_i = M_{i-1}$ for $1 \leq i \leq m+2$. We~have
\[
M_{\xi(\Filt)}=\mathop{\bigoplus}_{i=1}^{m+1} \frac{Y_{m+2-i}^\perp}{M_{m+2-i}^\perp}.
\]
Since $(\Gr_{\xi(\Filt)}(V), \theta_2)$ is not semistable and $M_{\xi(\Filt)}$ is the maximal destabilizer, we~have
\[
\deg(M_{\xi(\Filt)})>0.
\]
However, Lemma \ref{preserve} tells us that $\xi(\Filt)$ is quasi \grsemistable, and then
\begin{equation*}
\deg(M_{\xi(\Filt)})=\sum_{0\leq i\leq m+1}\big(\deg(Y_i)-\deg(M_i)\big)\leq 0
\end{equation*}
by Condition \eqref{cond:I} in Definition \ref{defn_gr_semi_odd}. This is a contradiction.
\end{enumerate}

The above discussion shows that $\mu(M_{\xi(\Filt)})= \mu(M_{\Filt})$ only if when $\Phi$ is injective. Then the second statement also holds. Moreover, if the equality
\[
\big(\mu(M_{\xi(\Filt)},\ \rk(M_{\xi(\Filt)})\big)= \big(\mu(M_{\Filt}),\rk(M_{\Filt})\big)
\]
holds, clearly $\Phi$ is an isomorphism on an open subset $U \subset X$ with $\codim(X/U) \geq 2$, which gives the third statement.
\end{proof}

\begin{lem}\label{lem_const}
Let $(V, \langle \,, \, \rangle, \nabla)$ be semistable and quasi \grsemistable. If the sequence of pairs $\{\big(\mu(M_{\xi^k(\Filt)}),\ \rk(M_{\xi^k(\Filt)})\big)\}_{k\geq 0}$ is constant, then the graded orthogonal Higgs sheaf $(\Gr_{\Filt}(V), \overline{\langle \,, \, \rangle}, \theta)$ is semistable.
\end{lem}

\begin{proof}
By Proposition \ref{prop_HN_fil}, it is equivalent to work on $(\Gr_{\Filt}(V), \theta)$ and we assume that $(\Gr_{\Filt}(V), \theta)$ is not semistable. We~use the same notations as in the proof of Lemma \ref{decrease}:
\begin{gather*}
\Filt: 0= L_0 \subsetneq L_1 \subsetneq \cdots \subsetneq L_m \subsetneq L_m^\perp\subsetneq\cdots \subsetneq L_1^\perp\subsetneq L_0^\perp= V,
\\
\Gr_{\Filt}(V)=\mathop{\bigoplus}_{i=1}^{m} \frac{L_i}{L_{i-1}}\oplus \frac{L_m^\perp}{L_m}\oplus \mathop{\bigoplus}_{i=1}^{m} \frac{L_{m-i}^\perp}{L_{m-i+1}^\perp},
\\
M_{\Filt}=\mathop{\bigoplus}_{i=1}^{m} \frac{M_i}{L_{i-1}}\oplus \frac{M_{m+1}}{L_m}\oplus \mathop{\bigoplus}_{i=1}^{m} \frac{N_{m-i+1}^\perp}{L_{m-i+1}^\perp},
\\
\xi(\Filt):0=M_0\subset M_1\subset\cdots\subset M_{m+1}\subset M_{m+1}^\perp\subset\cdots\subset M_1^\perp\subset M_0^\perp=V,
\\
\Gr_{\xi(\Filt)}(V)=\mathop{\bigoplus}_{i=1}^{m+1} \frac{M_i}{M_{i-1}}\oplus \frac{M_{m+1}^\perp}{M_{m+1}}\oplus \mathop{\bigoplus}_{i=1}^{m+1} \frac{M_{m+1-i}^\perp}{M_{m+2-i}^\perp},
\\
M_{\xi(\Filt)}=\mathop{\bigoplus}_{i=1}^{m+1} \frac{X_i}{M_{i-1}}\oplus \frac{X_{m+2}}{M_{m+1}}\oplus \mathop{\bigoplus}_{i=1}^{m+1} \frac{Y_{m+2-i}^\perp}{M_{m+2-i}^\perp}.
\end{gather*}
By Lemma \ref{decrease}, the morphism $\Phi:M_{\xi(\Filt)}\to M_{\Filt}$ is an isomorphism over an open subset $U$, where $\codim(X/U) \geq 2$. Now we always restrict to $U$, and say $\Phi$ is an isomorphism for convenience. Now we suppose that $\xi^{m+2}(\Filt)$ is of the form
\[
\xi^{m+2}(\Filt): 0= Z_0 \subset Z_1 \subset \cdots \subset Z_{2m+2}\subset Z_{2m+2}^\perp\subset\cdots \subset Z_1^\perp\subset Z_0^\perp= V.
\]
The isomorphisms
\begin{align*}
M_{\xi^{m+2}(\Filt)} \to M_{\xi^{m+1}(\Filt)} \to \cdots \to M_{\Filt}
\end{align*}
imply
\[
M_{\xi^{m+2}(\Filt)}=Z_1\oplus \frac{Z_2}{Z_1}\oplus\cdots\oplus \frac{Z_{2m+1}}{Z_{2m}},
\]
and thus $Z_{2m+1}$ is a $\nabla$-invariant subsheaf of $V$. We~also want to remind the reader that $\nabla(Z_{2m+2}) = 0$. Since $M_{\xi^{m+2}(\Filt)}$ is the maximal destabilizer, we~have
\[
\mu(Z_{2m+1})=\mu(M_{\xi^{m+2}(\Filt)})>0.
\]
This contradicts with the condition that $(V,\langle \,, \, \rangle,\nabla)$ is semistable.
\end{proof}

\begin{proof}[Proof of Theorem \ref{main thm}]
First, we~take an arbitrary quasi \grsemistable filtration $\Filt$ of $(V,\langle \,, \, \rangle,\nabla)$.
The sequence $\{\big(\mu(M_{\xi^k(\Filt)}),\ \rk(M_{\xi^k(\Filt)})\big)\}_{k\geq 0}$ decreases in the lexicographic ordering as $k$ grows by Lemma \ref{decrease}.
Thus, there exists a positive integer $k_0$ such that the sequence
\[
\{\big(\mu(M_{\xi^k(\Filt)}),\ \rk(M_{\xi^k(\Filt)})\big)\}_{k\geq k_0}
\]
becomes constant.
By Lemma \ref{lem_const}, we~finish the proof of this theorem.
\end{proof}

\begin{rem}
The statement of Lemma \ref{decrease} also works for arbitrary orthogonal bundle $(V, \langle \,, \, \rangle, \nabla)$, which could be unstable, while Lemma \ref{lem_const} only holds for semistable triples $(V, \langle \,, \, \rangle, \nabla)$.
\end{rem}

Combining Proposition \ref{ssregular} with Theorem \ref{main thm}, we~have:

\begin{thm}\label{thm_gr_and_quasi_gr_odd}
Let $(V,\langle \,, \, \rangle,\nabla)$ be a semistable orthogonal sheaf together with an orthogonal connection $\nabla$ such that the rank of $V$ is odd. Then, $(V,\langle \,, \, \rangle,\nabla)$ is \grsemistable if and only if $(V,\langle \,, \, \rangle,\nabla)$ is quasi \grsemistable.
\end{thm}

\section{Quasi \texorpdfstring{\protect\grsemistable}{grsemistable} filtration: even rank}\label{sect_even}
In this section, we~discuss the quasi \grsemistable (orthogonal/symplectic) filtrations of $(V, \langle \,, \, \rangle, \nabla)$, where $\rk(V)$ is even. In~the even rank case, there are two types of orthogonal/symplectic Hodge filtrations, which depend on the parity of lengths of the filtrations. Despite this fact, the proof of the main result (Theorem \ref{thm_even}) in the even rank case is almost the same as that of Theorem \ref{main thm} in the odd rank case. Therefore, we~only give constructions and state results rather than giving a detailed proof.

\begin{defn}
An orthogonal/symplectic Hodge filtration
\[
\Filt: 0=\Fil_0\subsetneq \Fil_1 \subsetneq \Fil_2 \subsetneq \cdots\subsetneq \Fil_t=V
\]
of $(V, \langle \,, \, \rangle, \nabla)$ is an \emph{odd} (\resp \emph{even}) type filtration if its length $t$ is odd (\resp even).
\end{defn}

\skpt
\begin{rem}
\begin{enumerate}
\item An odd type orthogonal Hodge filtration $\Filt$ can be written as
\[
\Filt: 0= L_0 \subsetneq L_1 \subsetneq \cdots \subsetneq L_m \subsetneq L_m^\perp\subsetneq\cdots \subsetneq L_1^\perp\subsetneq L_0^\perp= V,
\]
where $L_m$ is not Lagrangian (\ie $L_m\neq L_m^\perp$) and satisfies $\nabla(L_m)\subseteq L_m^\perp \otimes\Omega_X(\log D)$.

\item An even type orthogonal Hodge filtration $\Filt$ can be regarded as
\[
\Filt: 0= L_0 \subsetneq L_1 \subsetneq \cdots\subsetneq L_{m-1} \subsetneq L_m \subsetneq L_{m-1}^\perp\subsetneq\cdots \subsetneq L_1^\perp\subsetneq L_0^\perp= V,
\]
where $L_m$ is Lagrangian (\ie $L_m= L_m^\perp$) subsheaf and satisfies
\[
\nabla(L_m)\subseteq L_{m-1}^\perp \otimes\Omega_X(\log D).
\]
\end{enumerate}
In summary, the odd type orthogonal Hodge filtrations with even rank are similar to orthogonal Hodge filtrations with odd rank.
\end{rem}

\skpt
\begin{defn}[Quasi \grsemistable filtration in even rank case]\label{defn_gr_semi_even}
Given $(V,\langle \,, \, \rangle,\nabla)$, let $\Filt$ be an orthogonal/symplectic Hodge filtration. 
\begin{enumerate}
\item If
\[
\Filt: 0= L_0 \subsetneq L_1 \subsetneq \cdots \subsetneq L_m \subsetneq L_m^\perp\subsetneq\cdots \subsetneq L_1^\perp\subsetneq L_0^\perp= V
\]
is an odd type filtration, $\Filt$ is \emph{quasi \grsemistable} if the following two conditions hold: 
\begin{enumerate}\renewcommand{\theenumii}{\Roman{enumii}}
\item\label{cond2:I}
For any saturated Griffiths transverse filtration
\[
0=W_0\subset W_1 \subset \cdots\subset W_m
\]
such that $L_{i-1}\subset W_i\subset L_i$, we~have
\[
\sum_{0\leq i\leq m}\deg(W_i)\leqslant\sum_{0\leq i\leq m}\deg(L_i).
\]

\item\label{cond2:II}
For any saturated Griffiths transverse flitration
\[
0\subset S_1 \subset \cdots\subset S_m\subset S_{m+1}\subset T_m^\perp\subset\cdots \subset T_1^\perp\subset V
\]
such that
\[
\deg(S_{m+1})>\sum_{0\leq i\leq m}\big(\deg(L_i)-\deg(S_i)\big)+\sum_{0\leq i\leq m}\big(\deg(L_i)-\deg(T_i)\big),
\]
we have
\[
\nabla(S_{m+1})\subseteq S_{m+1}^\perp \otimes\Omega_X(\log D),
\]
where $L_{i-1}\subset S_i \subset T_i \subset L_i,\ L_m\subset S_{m+1}\subset L_m^\perp$ and $S_{m+1}$
is isotropic.
\end{enumerate}

\item If
\[
\Filt: 0= L_0 \subsetneq L_1 \subsetneq \cdots\subsetneq L_{m-1} \subsetneq L_m \subsetneq L_{m-1}^\perp\subsetneq\cdots \subsetneq L_1^\perp\subsetneq L_0^\perp= V
\]
is an even type filtration, $\Filt$ is \emph{quasi \grsemistable} if the following two conditions hold:
\begin{enumerate}\renewcommand{\theenumi}{\Roman{enumi}}
\item\label{enum:I}
For any saturated Griffiths transverse filtration
\[
0=W_0\subset W_1 \subset \cdots\subset W_m
\]
such that $L_{i-1}\subset W_i\subset L_i$, we~have
\[
\sum_{0\leq i\leq m}\deg(W_i)\leqslant\sum_{0\leq i\leq m}\deg(L_i).
\]

\item\label{enum:II}
For any saturated Griffiths transverse flitration
\[
0\subset S_1 \subset \cdots\subset S_m\subset T_m^\perp\subset\cdots \subset T_1^\perp\subset V
\]
such that
\[
\deg(S_m)>\sum_{0\leq i\leq m-1}\big(\deg(L_i)-\deg(S_i)\big)+\sum_{0\leq i\leq m}\big(\deg(L_i)-\deg(T_i)\big),
\]
we have
\[
\nabla(S_m)\subseteq L_m \otimes\Omega_X(\log D),
\]
where $L_{i-1}\subset S_i \subset T_i \subset L_i$ for $1\leq i\leq m$.
\end{enumerate}
\end{enumerate}
Furthermore, we~say that $(V,\langle \,, \, \rangle,\nabla)$ is \emph{quasi \grsemistable} (or $\nabla$ is a \emph{quasi \grsemistable \lconnection}) if there exists a quasi \grsemistable filtration on $V$.
\end{defn}

The following two propositions can be proved in the same manner as Propositions~\ref{ssregular} and \ref{positive} in the odd rank case.

\begin{prop}\label{prop_ssregular_even}
If an orthogonal/symplectic filtration $\Filt$ of $(V,\langle \,, \, \rangle,\nabla)$ is \grsemistable, then it is also quasi \grsemistable.
\end{prop}

\begin{prop}\label{positiveeven}
Let
\[
\Filt: 0= L_0 \subsetneq L_1 \subsetneq \cdots \subsetneq L_t = V
\]
be a nontrivial quasi \grsemistable filtration (either of odd or even type) of $(V,\langle \,, \, \rangle,\nabla)$. Then we have $\deg(L_i)\geq 0$. Moreover, $\deg(F)\leq \deg(L_1)$ holds for all subbundles $F\subset L_1$.
\end{prop}

As we did in Section~\ref{sect_odd}, we~have an operator $\xi$ on the set of quasi \grsemistable filtrations.
\begin{enumerate}
\item If
\[
\Filt: 0= L_0 \subsetneq L_1 \subsetneq \cdots \subsetneq L_m \subsetneq L_m^\perp\subsetneq\cdots \subsetneq L_1^\perp\subsetneq L_0^\perp= V
\]
is of odd type, we~define the orthogonal Higgs sheaf $(\Gr_{\Filt}(V), \overline{\langle \,, \, \rangle},\theta)$, where
\[
\Gr_{\Filt}(V)=\mathop{\bigoplus}_{i=1}^{m} \frac{L_i}{L_{i-1}}\oplus \frac{L_m^\perp}{L_m}\oplus \mathop{\bigoplus}_{i=1}^{m} \frac{L_{m-i}^\perp}{L_{m-i+1}^\perp}.
\]
If $(\Gr_{\Filt}(V),\theta)$ is not semistable, we~have a maximal destabilizer $M_{\Filt}$ and by Lem\-ma~\ref{destabilizer}, we~have
\[
M_{\Filt}=\mathop{\bigoplus}_{i=1}^{m} \frac{M_i}{L_{i-1}}\oplus \frac{M_{m+1}}{L_m}\oplus \mathop{\bigoplus}_{i=1}^{m} \frac{N_{m-i+1}^\perp}{L_{m-i+1}^\perp},
\]
where $L_{i-1}\subset M_i\subset N_i\subset L_i$ for $1\leq i\leq m$ and $M_{m+1}$ is isotropic. Based on the maximal destablizer, we~construct a new filtration $\xi(\Filt)$ as follows:
\[
\xi(\Filt):0=M_0\subset M_1\subset\cdots\subset M_{m+1}\subset M_{m+1}^\perp\subset\cdots\subset M_1^\perp\subset M_0^\perp=V.
\]

\item If
\[
\Filt: 0= L_0 \subsetneq L_1 \subsetneq \cdots\subsetneq L_{m-1} \subsetneq L_m \subsetneq L_{m-1}^\perp\subsetneq\cdots \subsetneq L_1^\perp\subsetneq L_0^\perp= V
\]
is of even type, we~define $(\Gr_{\Filt}(V),\langle \ \rangle,\theta)$ in a similar way, where
\[
\Gr_{\Filt}(V)=\mathop{\bigoplus}_{i=1}^{m} \frac{L_i}{L_{i-1}}\oplus \mathop{\bigoplus}_{i=1}^{m} \frac{L_{m-i}^\perp}{L_{m-i+1}^\perp}.
\]
If $(Gr_{\Filt}(V),\theta)$ is not semistable, we~have the maximal destabilizer
\[
M_{\Filt}=\mathop{\bigoplus}_{i=1}^{m} \frac{M_i}{L_{i-1}}\oplus \mathop{\bigoplus}_{i=1}^{m} \frac{N_{m-i+1}^\perp}{L_{m-i+1}^\perp},
\]
where $L_{i-1}\subset M_i\subset N_i\subset L_i$ for $1\leq i\leq m$. Then, the new filtration $\xi(\Filt)$ is defined~as
\[
\xi(\Filt):0=M_0\subset M_1\subset\cdots\subset M_m \subset M_{m+1}\subset M_m^\perp\subset\cdots\subset M_1^\perp\subset M_0^\perp=V,
\]
where $M_{m+1}:=L_m$.
\end{enumerate}

\begin{lem}
Let $\Filt$ be a quasi \grsemistable filtration of odd type (\resp even type), then $\xi(\Filt)$ is a Hodge filtration of odd type (\resp even type). Moreover, $\xi(\Filt)$ is also quasi \grsemistable of odd type (\resp even type).
\end{lem}

\begin{proof}
The proof of this lemma is similar to Lemma \ref{preserve}.
\end{proof}

\begin{lem}
Let $\Filt$ be a quasi \grsemistable filtration of $(V,\langle \,, \, \rangle,\nabla)$ in the even rank case. Assume that both $(\Gr_{\Filt}(V),\theta_1)$ and $(\Gr_{\xi(\Filt)}(V),\theta_2)$ are not semistable. Let $M_{\Filt}$ (\resp $M_{\xi(\Filt)})$ be the maximal destabilizer of $(\Gr_{\Filt}(V),\theta_1)$ (\resp $(\Gr_{\xi(\Filt)}(V),\theta_2))$. Then we have a natural morphism
\[
\Phi: M_{\xi(\Filt)}\to M_{\Filt}
\]
such that
\begin{enumerate}
\item $\mu(M_{\xi(\Filt)})\leq \mu(M_{\Filt})$,

\item if $\mu(M_{\xi(\Filt)})= \mu(M_{\Filt})$, the morphism $\Phi$ is injective, and then we have $\rk(M_{\xi(\Filt)})\leq\rk(M_{\Filt})$,

\item if $\big(\mu(M_{\xi(\Filt)},\ \rk(M_{\xi(\Filt)})\big)= \big(\mu(M_{\Filt}),\rk(M_{\Filt})\big)$, then $\Phi$ is an isomorphism on an open subset $U \subset X$ with $\codim(X/U) \geq 2$.
\end{enumerate}
\end{lem}

\begin{proof}
This lemma is an analogue of Lemma \ref{decrease}.
\end{proof}

\begin{thm}\label{main thm even}
Let $(V,\langle \,, \, \rangle,\nabla)$ be semistable and quasi \grsemistable with even rank. There exists a \grsemistable filtration on $(V,\langle \,, \, \rangle,\nabla)$.
\end{thm}

\begin{proof}
This theorem is an analogue of Theorem \ref{main thm}.
\end{proof}

Combining Proposition \ref{prop_ssregular_even} with Theorem \ref{main thm even}, we~have the main result in the even rank case.
\begin{thm}\label{thm_even}
Let $(V,\langle \,, \, \rangle,\nabla)$ be a semistable orthogonal/symplectic sheaf of even rank together with a \lconnection. Then $(V,\langle \,, \, \rangle,\nabla)$ is \grsemistable if and only if $(V,\langle \,, \, \rangle,\nabla)$ is quasi \grsemistable.
\end{thm}


\section{Classification of quasi \texorpdfstring{\protect\grsemistable}{grsemistable} orthogonal \texorpdfstring{\protect\lconnections}{lconnections} in small ranks}\label{sect_small_rank}

In this section, $(V, \langle \,, \, \rangle, \nabla)$ is an orthogonal sheaf together with an orthogonal \lconnection $\nabla$. We~make a careful discussion on the existence of quasi \grsemistable filtrations when $\rk (V) \leq 6$. We~prove that $(V, \langle \,, \, \rangle, \nabla)$ is always quasi \grsemistable when $\rk(V) \leq 4$ (Proposition \ref{rank4}), and we classify quasi \grsemistable \lconnections when $\rk(V) = 5 \text{ and }6$ (Proposition \ref{irregular}, \ref{classfication5} and \ref{classfication6}). Moreover, Proposition \ref{irregular} gives a construction of $(V, \langle \,, \, \rangle, \nabla)$ which is not quasi \grsemistable when $\rk(V) \geq 5$.

\subsection{Case I: ${\rk(V) \leq 4}$}
We start with a discussion of rank one isotropic saturated subsheaves $N\subset V$. Note that if a global section of $N$ vanishes on an open subset~$U$, then it is trivial. For this reason, it is enough to work on an open subset $U \subset X$, on~which $N|_U$ is locally free. Let $e \in N(U)$ be a local section. We~have
\[
\langle \nabla(e),e \rangle+\langle e,\nabla(e) \rangle=\lambda d\langle e,e \rangle=0,
\]
which implies $\langle \nabla(e),e \rangle=0$. Therefore,
\[
\nabla(N)\subset N^\perp\otimes\Omega_X(\log D)
\]
for any isotropic saturated subsheaf $N$ of rank one. Moreover, this property does not hold in the symplectic case, \ie $\langle \,, \, \rangle$ is skew-symmetric. Based on this fact, we~have the following lemma about a special orthogonal filtration
\[
0\subset N\subset N^\perp\subset V.
\] 

\begin{lem}\label{criterion}
Given $(V,\langle \,, \, \rangle,\nabla)$, let $N$ be a rank one isotropic saturated subsheaf. The filtration
\[
0\subset N\subset N^\perp\subset V
\]
is quasi \grsemistable if and only if the following two conditions hold:
\begin{enumerate}
\item $\deg(N)\geq 0$,

\item if an isotropic saturated subsheaf $W$ containing N satisfies one of the following conditions:
\begin{enumerate}
\item $\deg(W)>2\deg(N);$

\item $\deg(N)<\deg(W)\leq2\deg(N)$, and $\nabla(N)\subset W^\perp\otimes \Omega_X(\log D);$

\item $0<\deg(W)\leq \deg(N)$, and $\nabla(N)\subset W\otimes \Omega_X(\log D)$,
\end{enumerate}
then we must have $\nabla(W)\subset W^\perp\otimes \Omega_X(\log D)$.
\end{enumerate}
\end{lem}

\begin{proof} 
This is a direct consequence of Example \ref{smallm}.
\end{proof}

\begin{lem}\label{cork1}
Given $(V,\langle \,, \, \rangle,\nabla)$, let $N \subset W \subset V$ be isotropic saturated subsheaves of $V$ such that $\rk(N)=\rk(W)-1$. Then, $\nabla(W)\subset W^\perp\otimes \Omega_X(\log D)$ if and only if $\nabla(N)\subset W^\perp\otimes \Omega_X(\log D)$.
\end{lem}
\begin{proof}
One direction is clear, and we only have to prove that if
\[
\nabla(N)\subset W^\perp\otimes \Omega_X(\log D),
\]
we~have $\nabla(W)\subset W^\perp\otimes \Omega_X(\log D)$. Since this is a local property, we~work on an open subset $U \subset X$, on which $N$ is locally free. For convenience, we~directly suppose that
\[
N=\mathcal{O}_X\{e_1, e_2,\dots e_n\},\quad W=\mathcal{O}_X\{e_1, e_2,\dots e_n, e_{n+1}\}.
\]
Furthermore, $\nabla(W)\subset W^\perp\otimes \Omega_X(\log D)$ if and only if
\[
\langle \nabla(W),W \rangle=0
\]
holds. Therefore, it is equivalent to prove that $\langle \nabla(e_i), e_j \rangle=0$ holds for all $i,j$ such that $1\leq i,j\leq n+1$. Since $\nabla(N)\subset W^\perp\otimes \Omega_X(\log D)$, then $\langle \nabla(e_i), e_j \rangle=0$ holds for all $i$ such that $1\leq i\leq n$ and $1\leq j\leq n+1$. Thus, we~only have to check $\langle \nabla(e_{n+1}), e_i \rangle=0$ for $1\leq i\leq n+1$. The equality
\[
\langle \nabla(e_{n+1}), e_{n+1} \rangle+\langle e_{n+1}, \nabla(e_{n+1}) \rangle=\lambda d\langle e_{n+1},e_{n+1} \rangle=0
\]
implies that $\langle \nabla(e_{n+1}), e_i \rangle=0$ holds for $i= n+1$, and the equality
\[
\langle \nabla(e_{n+1}), e_i \rangle+\langle e_{n+1}, \nabla(e_i) \rangle=\lambda d\langle e_{n+1},e_i \rangle=0
\]
and the condition $\langle e_{n+1}, \nabla(e_i) \rangle=0$ implies that $\langle \nabla(e_{n+1}), e_i \rangle=0$ holds for all $i$ such that $1\leq i\leq n$. This finishes the proof of this lemma.
\end{proof}

Lemma \ref{criterion} and Lemma \ref{cork1} give the following result:
\begin{lem}\label{criterionrk5}
Suppose that $\rk(V) \leq 5$. Let $N$ be a rank one isotropic saturated subsheaf. Then the filtration
\[
0\subset N\subset N^\perp\subset V
\]
is a quasi \grsemistable filtration of $(V, \langle \,, \, \rangle, \nabla)$ if and only if the following two conditions hold:
\begin{enumerate}
\item $\deg(N)\geq 0$,

\item if $W$ is a rank two isotropic saturated subsheaf such that $N \subset W$ and $\deg(W)>2\deg(N)$, then we have $\nabla(W)\subset W^\perp\otimes \Omega_X(\log D)$.
\end{enumerate}
\end{lem}

\begin{prop}\label{rank4}
Any $(V, \langle \,, \, \rangle, \nabla)$ with $\rk(V) \leq 4$ is quasi \grsemistable.
\end{prop}

\begin{proof}
It is trivial when $\rk(V)\leq3$, and then we only discuss the case $\rk(V)=4$. If $(V, \langle \,, \, \rangle)$ is a semistable orthogonal sheaf, then the trivial filtration\vspace*{-3pt}
\[
0\subset V
\]
is quasi \grsemistable.

Now we assume that $(V, \langle \,, \, \rangle)$ is unstable and let $M$ be the maximal destabilizer. If $M$ is of rank one, the filtration\vspace*{-3pt}
\[
0\subset M\subset M^\perp\subset V
\]
is quasi \grsemistable by Lemma \ref{criterionrk5} because there is no rank two isotropic subsheaf~$W$ such that $\deg(W) > 2 \deg(M)$. Then, we~suppose that $M$ is of rank two, \ie it is Lagrangian. If there exists a rank one subsheaf $N$ such that
\begin{itemize}
\item $\deg(N) \geq 0$ and $N \subset M$,
\item $\nabla(N) \subset M\otimes \Omega_X(\log D)$,
\end{itemize}
then the filtration\vspace*{-3pt}
\[
0\subset N\subset M\subset N^\perp\subset V
\]
is an even type quasi \grsemistable filtration, otherwise the filtration\vspace*{-3pt}
\[
0\subset M\subset V
\]
is an even type quasi \grsemistable filtration. In~conclusion, we~find a quasi \grsemistable filtration for all cases. This finishes the proof.
\end{proof}

\subsection{Case II: ${\rk(V) = 5}$}\label{subsect_rak_5}
Before we consider the case $\rk(V)=5$, we~first construct a special orthogonal sheaf $(V, \langle \,, \, \rangle)$ together with an orthogonal \lconnection $\nabla$, which is not quasi \grsemistable. We~start with a well-known lemma (see \cite[Lem.\,1.3.3]{HL97} for instance).

\begin{lem}\label{lem_max_min}
Let $M, N$ be torsion free sheaves. If $\mu_{\min}(M) > \mu_{\max}(N)$, then $\Hom(M,N) = 0$.
\end{lem}

\begin{prop}\label{irregular}
Let $(V, \langle \,, \, \rangle)$ be an unstable orthogonal sheaf of $\rk(V)\geq 5$. Suppose that
\begin{enumerate}
\item the maximal destabilizer $M$ is of rank two and stable,
\item $\frac{M^\perp}{M}$ is a stable orthogonal sheaf,
\item there exists a \lconnection $\nabla$ satisfying\vspace*{-3pt}
\[
\nabla(M)\not \subset M^\perp\otimes\Omega_X(\log D).
\]
\end{enumerate}
Then $(V,\langle \,, \, \rangle,\nabla)$ is not quasi \grsemistable. Moreover, $(V, \langle \,, \, \rangle, \nabla)$ is semistable.
\end{prop}

\begin{proof}
We prove the first statement by contradiction, and assume that the orthogonal connection $(V,\langle \,, \, \rangle,\nabla)$ is quasi \grsemistable. Let
\[
\Filt: 0= L_0 \subsetneq L_1 \subsetneq L_2 \cdots \subsetneq L_1^\perp\subsetneq L_0^\perp= V
\]
be a quasi \grsemistable filtration.

Since $\sfrac{V}{M^\perp}\cong M^\vee$ is a stable bundle of $\deg(\sfrac{V}{M^\perp})<0$ and $\mu_{\min}(L_1)\geq 0$ by Proposition \ref{positive}, the composition
\[
L_1\hto V\onto \frac{V}{M^\perp}
\]
is zero by Lemma \ref{lem_max_min}. That is $L_1\subset M^\perp$.

Now we consider the composition
\[
L_1\hto M^\perp\onto \frac{M^\perp}{M}.
\]
Note that this composition is not surjective since $L_1$ is isotropic while $\sfrac{M^\perp}{M}$ is not isotropic. Then this map is zero because $\mu_{\min}(L_1)\geq 0$ and $\sfrac{M^\perp}{M}$ is a stable orthogonal sheaf of degree $0$ by assumption. This implies that $L_1\subset M$.

If $L_1=M$, then we have
\[
\nabla(M)\subset L_2\otimes \Omega_X(\log D)\subset M^\perp\otimes \Omega_X(\log D),
\]
which contradicts the assumption that $\nabla(M)\not \subset M^\perp\otimes\Omega_X(\log D)$. Therefore, $L_1$ is a proper subsheaf of $M$.

Now consider the natural morphism
\[
h: L_2\to \frac{L_1^\perp}{M^\perp}
\]
based on the above argument that $L_1$ is a proper subsheaf of $M$. If $h=0$, then $L_2\subset M^\perp$, and we have
\[
\nabla(L_1)\subset L_2\otimes \Omega_X(\log D)\subset M^\perp\otimes \Omega_X(\log D).
\]
This actually implies
\[
\nabla(M) \subset M^\perp\otimes\Omega_X(\log D)
\]
by Lemma \ref{cork1}, and we get a contradiction again. On the other hand, if $h\neq 0$, then
\[
\deg(\ker(h))\geq \deg(L_2)+\deg(M)-\deg(L_1)> \deg(L_2)+\deg(L_1).
\]
This inequality contradicts Condition \eqref{cond:I} in Definition \ref{defn_gr_semi_odd}. In~conclusion, $(V, \langle \,, \, \rangle, \nabla)$ is not quasi \grsemistable.

For the second statement, we~still prove it by contradiction. Suppose that $(V, \langle \,, \, \rangle, \nabla)$ is unstable. Denote by $W$ the maximal destabilizer of $(V, \langle \,, \, \rangle, \nabla)$. Clearly, $W$ is isotropic of $\deg(W)>0$ and $\nabla(W)\subset W\otimes \Omega_X(\log D)$.

If $W$ is a rank one saturated subsheaf, we~apply the same argument above for compositions
\begin{align*}
W \hto V \onto \frac{V}{M^{\perp}}, \quad W \hto M^\perp\onto \frac{M^\perp}{M}
\end{align*}
and obtain $W \subset M$. Therefore,
\[
\nabla(M)\subset M^\perp\otimes\Omega_X(\log D)
\]
by Lemma \ref{cork1}, which contradicts the assumption in the statement of the proposition. If $\rk(W)\geq 2$, we~have $W= M$. Since $W$ is $\nabla$-invariant, $M$ is also $\nabla$-invariant. This also gives a contradiction.

Finally, suppose that $W$ is unstable as an orthogonal sheaf. Denote by $N\subset W$ be the maximal destabilizer. With the same argument as above, we~have $N\subset M$ and $W\subset M^\perp$. Therefore,
\[
\nabla(N)\subset W\otimes \Omega_X(\log D)\subset M^\perp\otimes \Omega_X(\log D),
\]
and this implies that $\nabla(M) \subset M^\perp\otimes\Omega_X(\log D)$ by Lemma \ref{cork1}, which is a contradiction.

In conclusion, $(V, \langle \,, \, \rangle, \nabla)$ is semistable.
\end{proof}

In the following examples, we~use the criterion in Proposition \ref{irregular} and construct examples of semistable orthogonal connections, which are not \grsemistable, in both characteristic zero and characteristic $p$.

\begin{examp}\label{exmp_char_zero}
In this example, we~will construct an orthogonal connection of rank~$5$ in characteristic zero, which is not \grsemistable. This gives a counterexample for Simpson's question (Question \ref{quest_simp}) in orthogonal group case.

Let $X$ be a complex smooth projective curve with genus $g \geq 2$. We~fix a rank $2$ stable vector bundle $M$ on $X$ with $\deg(M)=d>0$, and we suppose $g > d+1$. Then we choose a nontrivial extension
\[
0 \to \mathcal{O}_X \to M' \to \det(M^{\vee}) \to 0.
\]
Note that such a nontrivial extension always exists because
\[
\dim \Ext^1(\det(M^{\vee}), \mathcal{O}_X ) \geq g-d-1 > 0
\]
by Riemann-Roch theorem. Now we define $E := M \otimes M'$ and there is a natural orthogonal structure given as follows
\begin{align*}
\langle \,, \, \rangle_E : E \otimes E & \To{\cong} M \otimes M \otimes M' \otimes M' \to \wedge^2 M \otimes \wedge^2 M' \\
& \To{\cong} \det(M) \otimes \det(M^{\vee}) \cong \mathcal{O}_X.
\end{align*}
Moreover, Choe and Hitching showed that any rank $4$ orthogonal bundle, which contains rank $2$ isotropic subbundle, is of the above form \cite[Th.\,4.2]{CH23}. Then we obtain a rank $4$ orthogonal bundle $(E, \langle \,, \, \rangle_E)$. We~\emph{claim} that there exists an orthogonal connection $\nabla_0$ on $E$. Now we are ready to define a rank $5$ orthogonal connection $(V, \langle \,, \, \rangle, \nabla)$, where
\begin{itemize}
\item $V:= E \oplus \mathcal{O}_X$,
\item the orthogonal structure $\langle \,, \, \rangle$ on $V$ is induced by $\langle \,, \, \rangle_E$ and the trivial one on $\mathcal{O}_X$,
\item $\nabla:=\nabla_0 \oplus \mathrm{d}_\mathrm{can}$, where $\mathrm{d}_\mathrm{can}: \mathcal{O}_X \to \Omega_X$ is the exterior differential.
\end{itemize}
For this orthogonal connection $(V, \langle \,, \, \rangle, \nabla)$, it has the following properties that the maximal destabilizer of $V$ is $M$, which is stable of rank $2$ with $\deg (M) =d > 0$, and $\nabla(M) \nsubseteq M^{\perp} \otimes \Omega_X$. Therefore, $(V, \langle \,, \, \rangle, \nabla)$ is not quasi \grsemistable by Proposition~\ref{irregular}. The claim $\nabla(M) \nsubseteq M^{\perp} \otimes \Omega_X$ can be seen as follows. Assume the contrary. As~$\nabla(M)=\nabla_0(M)\subset (M\otimes M^\prime)\otimes \Omega_X$, it follows that
\[
\nabla(M)\subset \big((M\otimes M^\prime)\cap M^{\perp}\big)\otimes \Omega_X=M\otimes \Omega_X.
\]
The last equality holds because $M^{\perp}=M\otimes \mathcal{O}_X\oplus \mathcal{O}_X=M\oplus \mathcal{O}_X$. Then we have $\deg(M)=0$ since $M$ admits a connection $\nabla|_{M}$, contradicting $\deg(M)>0$.

Now we will prove the claim that there exists an orthogonal connection $\nabla_0$ on $E = M \otimes M'$. By \cite[Th.\,4.2]{BMW11}, if there is no orthogonal connection on $E$, there exists a decomposition
\[
E= (W \oplus W^{\vee}) \oplus T, \quad \langle \,, \, \rangle_E = \langle \,, \, \rangle_W \oplus \langle \,, \, \rangle_T
\]
such that $\deg W > 0$. Then there are two possibilities for the isotropic subbundle $W$: $\rk(W)=1 \text{ or } 2$.
\begin{itemize}
\item Suppose $\rk(W)=1$. The short exact sequence
\[
0 \to \mathcal{O}_X \to M' \to \det(M^{\vee}) \to 0
\]
gives the following one
\[
0 \to M \to M \otimes M' \to M^{\vee} \to 0
\]
by tensoring with $M$. Then the composition
\[
W \hto M \otimes M' = E \to M^{\vee}
\]
is zero because $\deg(W) > \deg(M^{\vee})$. Thus $W \subseteq M$. With a similar argument, we~have $W^{\vee} \subseteq M$. Therefore, $W \oplus W^{\vee} \subseteq M$. As
\[
\frac{M}{W\oplus W^{\vee}} \subset \frac{E}{W\oplus W^{\vee}}\cong T,
\]
it follows that $ \spfrac{M}{W\oplus W^{\vee}}$ is torsion free. For reason of rank, it must be zero. In~other words, $W \oplus W^{\vee} = M$. However, this contradicts the assumption that $M$ is stable.

\item Suppose $\rk(W)=2$. In~this case,
\[
E = M \otimes M' = W \oplus W^\vee.
\]

\begin{enumerate}
\item[(a)] If $W$ is not semistable, let $L$ be the maximal destabilizer of $W$. Then
\[
\mu_{\min}(W) = \deg(W) - \deg(L) > - \deg L \geq - \frac{\deg(M)}{2} = \mu(M^{\vee}).
\]
Since $\mu_{\min}(W) > \mu(M^{\vee})$, the composition
\[
W \hto E \to M^{\vee}
\]
is zero by Lemma \ref{lem_max_min}. With the same argument as in the case of $\rk(W)=1$, we~have $W=M$. However, this contradicts the assumption that $M$ is stable.

\item[(b)] If $W$ is semistable, the argument in Case (a) implies that $W=M$. Therefore, we~have
\[
M \otimes M' = M \oplus M^{\vee}.
\]
Then the short exact sequence
\[
0 \to M \to M \otimes M' \to M^{\vee} \to 0
\]
splits. The following argument will show that the short exact sequence
\[
0 \to \mathcal{O}_X \to M' \to \det(M^{\vee}) \to 0
\]
splits, and thus we obtain a contradiction about the choice of $M'$ at the beginning of this example.

By Proposition \cite[Prop.\,4.5]{CH23}, any rank $2$ isotropic subbundle of $M \otimes M'$ is of the form $M \otimes L'$ or $L \otimes M'$, where $L \subseteq M$ and $L' \subseteq M'$ are subbundles. Since $E=W \oplus W^{\vee}$, we~have $W \cap W^{\vee} = 0$, which implies that $W^\vee = M \otimes L'$ for some line bundle $L' \subseteq M'$. We~have
\[
\deg(L') = -d = \deg(\det(M^{\vee})).
\]
Then the composition
\[
L' \hto M' \to \det(M^{\vee})
\]
is either zero or an isomorphism. If the composition is zero, then $L'$ is a subbundle of $\mathcal{O}_X$, which is a contradiction. Thus, the composition is an isomorphism, and we obtain a splitting $\det(M^{\vee}) \to M'$.
\end{enumerate}
\end{itemize}
In conclusion, neither of the above cases can happen. By \cite[Th.\,4.2]{BMW11}, there exists an orthogonal connection $\nabla_0$ on $E = M \otimes M'$.
\end{examp}

\begin{examp}\label{exmp_const}
In this example, we~will construct an orthogonal connection of rank~$5$ in positive characteristic, which is not quasi \grsemistable.

Let $X$ be a smooth projective curve over $k$ with genus $g$ such that
\begin{align*}
3\leq p < g-1,
\end{align*}
where $p$ is the characteristic of $k$. There exists a stable bundle $M'$ on $X$ such that $\rk(M') = 2$ and $\deg(M')=1$. Consider the Frobenius pullback $M:=F^*_X(M')$, of which $\rk(M)=2$ and $\deg(M)=p$. Although $M$ may not be semistable \cite{Gi73}, we~can choose a specific $M'$ such that $M$ is semistable. Since the degree $p$ and the rank $2$ are coprime, the semistable bundle $M$ is stable. Let $\nabla_0: M \to M \otimes \Omega_X$ be the canonical connection. Next, we~consider the sheaf $(\wedge^2 M^{\vee}) \otimes \Omega_X$. By Riemann--Roch formula, we~have
\begin{align*}
h^0((\wedge^2 M^\vee) \otimes \Omega_X) & \geq \deg( (\wedge^2 M^\vee) \otimes \Omega_X ) + (1-g) \\
& = (-p+2g-2) + (1-g) \\
& = -p + (g-1) > 0.
\end{align*}
Therefore, there exists a nontrivial section of $(\wedge^2 M^{\vee}) \otimes \Omega_X$. Let
\[
V:= M \oplus M^{\vee} \oplus \mathcal{O}_X
\]
be the rank $5$ orthogonal bundle with natural orthogonal structure $\langle \,, \, \rangle$. Then, the nontrivial section of $(\wedge^2 M^{\vee}) \otimes \Omega_X$ induces a nontrivial $\mathcal{O}_X$-linear morphism
\[
\theta_1: M \to M^{\vee} \otimes \Omega_X
\]
such that for local sections $a,b$ of $M$,
\[
\langle \theta_1(a),b \rangle_M + \langle a,\theta_1(b) \rangle_M = 0.
\]
By Proposition \ref{irregular}., the orthogonal connection
\begin{align*}
\nabla:=
\begin{pmatrix}
\nabla_0 & & \\
\theta_1 & -\nabla^\vee_0 & \\
& & \mathrm{d}_\mathrm{can}
\end{pmatrix}: V \to V \otimes \Omega_X,
\end{align*}
is not quasi \grsemistable, where $\mathrm{d}_\mathrm{can}: \mathcal{O}_X \to \Omega_X$ is the exterior differential. Under the equivalence of gr-semistability and quasi gr-semistability (Theorem \ref{thm_gr_and_quasi_gr_odd}), we~conclude that we find an example in rank $5$ which is not \grsemistable.
\end{examp}

\begin{examp}\label{exmp_idea}
In the above two examples, we~constructed semistable orthogonal connections which is not \grsemistable in the case of $D=\emptyset$. In this example, we~give a brief idea about constructing examples satisfying conditions given in Proposition~\ref{irregular} when $D\neq\emptyset$.

Let $X$ be a smooth projective curve with genus $g(X) \geq 2$. Let $M$ be a rank two stable bundle of positive degree, and there is a natural orthogonal structure on $M \oplus M^{\vee}$ and denote it by $\langle \,, \, \rangle_M$. Let $N$ be an arbitrary stable orthogonal bundle on~$X$ and the orthogonal structure of $N$ is denoted by $\langle \,, \, \rangle_N$. Define
\[
V=M\oplus M^\vee\oplus N.
\]
We have a natural orthogonal structure $\langle \,, \, \rangle$ on $V$ induced by $\langle \,, \, \rangle_M$ and $\langle \,, \, \rangle_N$.

Let $\nabla_0: M\to M\otimes\Omega_X(\log D)$ be a logarithmic \lconnection on $M$, and $\nabla_0^\vee: M^\vee \to M^\vee \otimes \Omega_X(\log D)$ its dual logarithmic connection of $\nabla_0$. Then,
\begin{align*}
\begin{pmatrix}
\nabla_0 & \\
& -\nabla^\vee_0
\end{pmatrix}: M \oplus M^\vee \to (M \oplus M^\vee) \otimes \Omega_X(\log D)
\end{align*}
is an orthogonal \lconnection on $M \oplus M^\vee$. We~also choose $\theta_1: M \to M^\vee \otimes \Omega_X(\log D)$ a nontrivial $\mathcal{O}_X$-linear morphism such that
\[
\langle \theta_1(a),b \rangle_M + \langle a,\theta_1(b) \rangle_M = 0
\]
for all local sections $a,b$ of $M$. We~may always find such a section by enlarging the support of $D$. Clearly,
\[
\begin{pmatrix}
\nabla_0 & \\
\theta_1 & -\nabla^\vee_0
\end{pmatrix}: M \oplus M^\vee \to (M \oplus M^\vee) \otimes \Omega_X(\log D)
\]
is an orthogonal \lconnection. Choosing a \lconnection $\nabla_N$ on $N$, we~get an orthogonal \lconnection
\[
\nabla:=
\begin{pmatrix}
\nabla_0 & & \\
\theta_1 & -\nabla^\vee_0 & \\
& & \nabla_N
\end{pmatrix}: V \to V \otimes \Omega_X(\log D)
\]
on $(V, \langle \,, \, \rangle)$. From the above construction, we~have
\[
\nabla(M)\not\subset M^\perp\otimes\Omega_X(\log D).
\]
Thus, we~obtain an orthogonal bundle $(V, \langle \,, \, \rangle)$ of rank larger than $5$ with an orthogonal \lconnection $\nabla$ satisfying all conditions in Proposition \ref{irregular}.
\end{examp}

Proposition \ref{irregular} gives the construction of a special class of $(V, \langle \,,\, \rangle, \nabla)$, which is not quasi \grsemistable. Furthermore, we~can show that if $(V, \langle \,,\, \rangle, \nabla)$ is not quasi \grsemistable and if $\rk(V)=5 \text{ and } 6$, then it must satisfy the conditions in Proposition \ref{irregular}. In~the next proposition, we~shall discuss the case $\rk(V)=5$ first.

\begin{prop}\label{classfication5}
Suppose that $\rk(V)=5$. A triple $(V,\langle \,, \, \rangle,\nabla)$ is not quasi \grsemistable if and only if $(V, \langle \,, \, \rangle)$ is an unstable orthogonal sheaf such that the maximal destabilizer $M$ is stable of rank two and satisfies
\[
\nabla(M)\not\subset M^\perp\otimes\Omega_X(\log D).
\]
\end{prop}

\begin{proof} 
One direction is given by Proposition \ref{irregular}. We~consider the other direction and assume that $(V,\langle \,, \, \rangle,\nabla)$ is not quasi \grsemistable. Clearly, $(V, \langle \,, \, \rangle)$ is unstable, otherwise the trivial filtration
\[
0\subset V
\]
is quasi \grsemistable. Now denote by $M$ the maximal destabilizer of $V$. If $M$ is not a stable rank two bundle, there exists saturated subsheaf $N \subset M$ of rank one such that $\deg(N)=\mu(M)$. Then, the filtration
\[
0\subset N\subset N^\perp\subset V
\]
is quasi \grsemistable by Lemma \ref{criterionrk5}. Therefore, $M$ is stable of rank two. Furthermore, we~have
\[
\nabla(M)\not\subset M^\perp\otimes\Omega_X(\log D),
\]
otherwise the filtration
\[
0\subset M \subset M^\perp\subset V
\]
is quasi \grsemistable.
\end{proof}

\subsection{Case III: ${\rk(V) = 6}$}
For the rank $6$ case, we~first prove two lemmas.

\begin{lem}\label{degreeincreasing}
Suppose that $(V,\langle \,, \, \rangle,\nabla)$ is not quasi \grsemistable of rank $6$. Let
\[
\Filt_{LW}:0\subset L\subset W\subset L^\perp\subset V
\]
be an orthogonal Hodge filtration such that
\begin{itemize}
\item $\rk(L)=1, \ \rk(W)=3$ and $\deg(L)>0, \ \deg(W)>0$,
\item $\deg(L)+\deg(W)>\mu_{\max}(V)$,
\item $2\deg(L)+\deg(W)>2\mu_{\max}(V)$.
\end{itemize}
Then there exists a rank three isotropic subsheaf $W^\prime$ containing $L$ and satisfying one of the following conditions:
\begin{enumerate}
\item $\deg(W^\prime)>\max\{2\deg(L),\deg(W)\}$;
\item $\deg(W^\prime)>\deg(W)$, and the filtration
\[
\Filt_{LW}:0\subset L\subset W^\prime\subset L^\perp\subset V
\]
is an orthogonal Hodge filtration;
\item $\deg(W^\prime)>2\deg(L)$, and there exists a rank two subbundle $P_2$ of $W$ such that
\begin{align*}
& 2\deg(P_2)>\deg(W)+2\deg(L), \\
& \nabla(P_2)\subset P_2^\perp\otimes\Omega_X(\log D), \\ & \nabla(P_2)\not\subset W\otimes\Omega_X(\log D).
\end{align*}
\end{enumerate}
\end{lem}
\begin{proof}
Since $(V,\langle \,, \, \rangle,\nabla)$ is not quasi \grsemistable, the filtration $\Filt_{LW}$ does not satisfy Condition \eqref{cond2:I} or \eqref{cond2:II} in Definition \ref{defn_gr_semi_even}. If $\Filt_{LW}$ does not satisfy Condition~\eqref{cond2:I}, there exists a rank two isotropic sheaf $S$ between $L$ and $W$ such that
\begin{align*}
\deg(S)&>\deg(L)+\deg(W),
\\\
\tag*{or}
\deg(S)&>\deg(W) \text{ and } \nabla(L)\subset S\otimes\Omega_X(\log D).
\end{align*}

For the \emph{first case} $\deg(S)>\deg(L)+\deg(W)$, we~consider
\[
\Filt_S:0\subset S\subset S^\perp\subset V,
\]
which is not quasi \grsemistable. Clearly, $\mu_{\min}(S)>0 $ because
\[
\deg(S)>\deg(L)+\deg(W)>\mu_{\max}(V).
\]
Then there exist isotropic subsheaves
\[
0\subset T\subset U\subset S\subset W^\prime 
\]
such that $\rk(U)\leq1,\ \rk(W^\prime)=3$ and
\[
\nabla(U )\subset W^\prime\otimes\Omega_X(\log D),\ \deg(W^\prime)>2\deg(S)-\deg(U)-\deg(T).
\]
Then,
\begin{equation*}
\begin{split}
\deg(W^\prime)&> 2\deg(S)-\deg(U)-\deg(T)\\
&>\deg(W)+2\deg(L)+\deg(W)-\deg(U)-\deg(T)\\
&>\deg(W)+2\mu_{\max}(V)-\mu_{\max}(V)-\mu_{\max}(V)=\deg(W).
\end{split}
\end{equation*}
Moreover, $U=L$ or $\deg(U)\leq \deg(S)-\deg(L)$. If $U=L$, it is clear
\[
\nabla(L)\subset W^\prime\otimes\Omega_X(\log D).
\]
If $\deg(U)\leq \deg(S)-\deg(L)$, we have
\[
\deg(W^\prime)>2\deg(S)-2(\deg(S)-\deg(L))=2\deg(L).
\]

For the \emph{second case} that $\deg(S)>\deg(W)$ and $\nabla(L)\subset S\otimes\Omega_X(\log D)$, consider the filtration
\[
\Filt_{LS}:0\subset L\subset S\subset S^\perp\subset L^\perp\subset V,
\]
and there exists a rank three isotropic subsheaf $W^\prime$ containing $L$ such that $\deg(W^\prime)>2\deg(S)$ and
\[
\nabla(L)\subset W^\prime\otimes\Omega_X(\log D).
\]
Therefore,
\[
\deg(W^\prime)>2\deg(S)>2\deg(W)>\deg(W).
\]

If $\Filt_{LW}$ does not satisfy Condition \eqref{cond2:II}, then there exists a Hodge filtration
\[
P_1\subset P_2\subset P_2^\perp\subset Q_1^\perp\subset V
\]
such that
\begin{gather*}
0\subset P_1\subset Q_1\subset L \subset P_2\subset W,
\\
\rk(P_2)=2,\quad 2\deg(P_2)>\deg(W)+2\deg(L)-\deg(P_1)-\deg(Q_1),
\\
\tag*{and}
\nabla(P_2)\not\subset W\otimes\Omega_X(\log D).
\end{gather*}
Note that $Q_1=0$ or $L$. If $Q_1=0$, then $P_1=0$. Thus,
\[
\deg(P_2)>\dfrac{1}{2}(\deg(W)+2\deg(L))>\mu_{\max}(V).
\]
With a similar argument as the proof in the \emph{first case}, such a subsheaf $W^\prime$ exists by considering the filtration
\[
\Filt_{P_2}:0\subset P_2\subset P_2^\perp\subset V.
\]
If $Q_1=L$, a similar argument as in the \emph{second case} will produce such a subsheaf $W^\prime$ by considering
\[
\Filt_{LP_2}:0\subset L\subset P_2\subset P_2^\perp\subset L^\perp\subset V.
\]
This finishes the proof of this lemma.
\end{proof}

\begin{lem}\label{lagrange}
Suppose that $(V,\langle \,, \, \rangle,\nabla)$ is not quasi \grsemistable and is of rank~$6$. Define
\[
\alpha(V):=\max_{M}\ \{\deg(M)\ |\ M \text{ is a rank three isotropic subsheaf of $V$}\}.
\]
Then $\alpha(V)<2\mu_{\max}(V)$.
\end{lem}

\begin{proof}
We still prove this lemma by contradiction, and assume that $\alpha(V)\geq2\mu_{\max}(V)$. Let $M$ be a rank three isotropic subsheaf of $V$ such that $\deg(M)=\alpha(V)$. Then we have $\mu_{\min}(M)\geq0$ since $\deg(M)\geq 2\mu_{\max}(V)$. This implies that the filtration
\[
\Filt_{M}:0\subset M\subset V,
\]
which is not quasi \grsemistable, satisfies Condition \eqref{cond2:I}. There exist saturated subsheaves $S\subset T\subset M$ such that
\begin{align*}
& \deg(T)+\deg(S)>\deg(M), \\
& \nabla(S)\subset T^\perp\otimes\Omega_X(\log D), \\
& \nabla(S)\not\subset M\otimes\Omega_X(\log D).
\end{align*}
Then we have $\rk(T)=2$, otherwise
\[
\deg(T)+\deg(S)\leq \mu_{\max}(V)+\mu_{\max}(V)=2\mu_{\max}(V)\leq \deg(M),
\]
which is a contradiction. In~summary, there exists a rank two subsheaf $T$ of $M$ such that
\[
\deg(T)>\frac{\deg(M)}{2}\geq \mu_{\max}(V),\ \nabla(T)\subset T^\perp\otimes\Omega_X(\log D),
\]
and
\[
\nabla(T)\not\subset M\otimes\Omega_X(\log D).
\]
Moreover, $\mu_{\min}(T)\geq 0$.

Now we choose such a subsheaf $T$ with maximal degree and consider the orthogonal Hodge filtration
\[
\Filt_{T}:0\subset T\subset T^\perp\subset V.
\]
Since it is not quasi \grsemistable, there exist saturated subsheaves $S_1\subset T_1\subset T\subset S_2$ such that $S_2$ is isotropic of rank three and
\[
\nabla(T_1)\subset S_2\otimes\Omega_X(\log D),\ \nabla(S_2)\not\subset S_2\otimes\Omega_X(\log D),
\]
and
\[
\deg(S_2)>2\deg(T)-\deg(T_1)-\deg(S_1).
\]
Since $\nabla(S_2)\not\subset S_2\otimes\Omega_X(\log D)$, we~have $\rk(T_1)\leq1$ by Lemma \ref{cork1}. If $T_1=0$, clearly
\[
\deg(S_2)>2\deg(T)\geq \deg(M).
\]
This contradicts our choice of $M$. Therefore, we~have $\rk(T_1)=1$. Now if $\rk(T_1)=1$, we~get the following inequality
\begin{equation*}
\begin{split}
\deg(\sfrac{S_2}{T})-\deg(\sfrac{T^\perp}{M})&=\deg(M)+\deg(S_2)-2\deg(T)\\
&>\deg(M)-\deg(T_1)-\deg(S_1)\\
&\ge2\mu_{\max}(V)-\mu_{\max}(V)-\mu_{\max}(V)=0,
\end{split}
\end{equation*}
which implies $S_2=M$ by the natural morphism $\sfrac{S_2}{T}\to \sfrac{T^\perp}{M}$. Moreover, we~have $\deg(T_1)>0$.

Then, we~consider the orthogonal Hodge filtration
\[
\Filt_{T_1M}:0\subset T_1\subset M\subset T_1^\perp \subset V.
\]
By Lemma \ref{degreeincreasing}, there exists a rank two saturated subsheaf $P_2 \subset M$ such that
\begin{align*}
& 2\deg(P_2)>\deg(M)+2\deg(T_1), \\
& \nabla(P_2)\subset P_2^\perp\otimes\Omega_X(\log D), \\
& \nabla(P_2)\not\subset M\otimes\Omega_X(\log D).
\end{align*}
Then,
\begin{equation*}
\begin{split}
2\deg(P_2)&>\deg(M)+2\deg(T_1)\\
&>2\deg(T)-\deg(T_1)-\deg(S_1)+2\deg(T_1)\\
&=2\deg(T)+\deg(T_1)-\deg(S_1)\\
&\geq 2\deg(T).
\end{split}.
\end{equation*}
This contradicts our choice of $\deg(T)$, which is of maximal degree. In~conclusion, we~have $\alpha(V)<2\mu_{\max}(V)$.
\end{proof}


\begin{prop}\label{classfication6}
Let $(V,\langle \,, \, \rangle,\nabla)$ be an orthogonal sheaf of rank~$6$ together with an orthogonal \lconnection. Then $(V,\langle \ \rangle,\nabla)$ is not quasi \grsemistable if and only if~$V$ is an unstable orthogonal sheaf such that
\begin{enumerate}
\item the maximal destabilizer $M$ is a stable rank two saturated subsheaf,
\item $\sfrac{M^\perp}{M}$ is stable or $\sfrac{M^\perp}{M}$ is a direct sum of two non isomorphic orthogonal torsion free sheaves of rank one,
\item $\nabla(M)\not\subset M^\perp\otimes\Omega_X(\log D)$.
\end{enumerate}
\end{prop}

\begin{proof}
One direction is exactly proved by Proposition \ref{irregular}. Now we assume that $(V,\langle \ \rangle,\nabla)$ is not quasi \grsemistable. Clearly, the trivial filtration $0 \subset V$ is not quasi \grsemistable, and then $V$ is unstable. We~prove the other three statements by contradiction.

Let $M$ be the maximal destabilizer of $V$. Lemma \ref{lagrange} shows that $\rk(M) \leq 2$. If $M$ is not stable of rank two, then there exists a saturated subsheaf $N \subset M$ of rank one such that $\deg(N)=\mu(M)=\mu_{\max}(V)$. Now consider the orthogonal Hodge filtration
\[
\Filt_{N}:0\subset N\subset N^\perp\subset V.
\]
Since it is not quasi \grsemistable, there exists a Hodge filtration
\[
R_1\subset R_2\subset U_1^\perp\subset V
\]
such that
\[
0\subset R_1\subset U_1\subset N \subset R_2,
\]
with $\rk(R_2)=2$ or $3$,
\[
\deg(R_2)>2\deg(N)-\deg(R_1)-\deg(U_1),
\]
and
\[
\nabla(R_2)\not\subset R_2^\perp\otimes\Omega_X(\log D).
\]
We first suppose $\rk(R_2)=2$. Then, the condition $\nabla(R_2)\not\subset R_2^\perp\otimes\Omega_X(\log D)$ and Lemma \ref{cork1} implies that $R_1 = U_1 = 0$. Then,
\[
\deg(R_2)>2\deg(N)=2\mu_{\max}(V),
\]
and this is a contradiction. Thus, $\rk(R_2)=3$. In~this case, $U_1=0$ or $U_1 = N$. With a similar argument as above, Lemma \ref{lagrange} implies $U_1=N$. In~summary, there exists a rank three isotropic subsheaf $R_2$ such that
\[
\deg(R_2)>0,\ \nabla(N)\subset R_2 \otimes \Omega_X(\log D),\ \nabla(R_2)\not\subset R_2 \otimes \Omega_X(\log D).
\]
Now we choose such a subsheaf $R_2$ with maximal degree and consider the orthogonal Hodge filtration
\[
\Filt_{NR_2}:0\subset N\subset R_2\subset N^\perp \subset V.
\]
By Lemma \ref{degreeincreasing}, there exists an isotropic subsheaf $W^\prime$ of rank three such that $\deg(W^\prime)>2\deg(N)=2\mu_{\max}(V)$. However, this contradicts Lemma \ref{lagrange}. Therefore, the maximal destabilizer $M$ of $V$ is stable of rank two.

For the second statement, suppose that $\frac{M^\perp}{M}$ is not stable. Then there exists a rank three isotropic subsheaf $Z$ contain $M$ such that $\deg(\frac{Z}{M})\geq 0$. We have $\deg(Z)\geq \deg(M)=2\mu_{\max}(V)$, which contradicts Lemma \ref{lagrange}. Therefore, $\frac{M^\perp}{M}$ is stable.

For the last statement, if $ \nabla(M)\subset M^\perp\otimes\Omega_X(\log D)$, we consider the Hodge filtration
\[
\Filt_{M}:0\subset M\subset M^\perp\subset V,
\]
which is not quasi \grsemistable. Then there exist saturated subsheaves $S_1\subset T_1\subset M\subset S_2$ such that $S_2$ is isotropic of rank three and
\[
\nabla(T_1)\subset S_2\otimes\Omega_X(\log D),\ \nabla(S_2)\not\subset S_2\otimes\Omega_X(\log D),
\]
and
\[
\deg(S_2)>2\deg(M)-\deg(T_1)-\deg(S_1).
\]
Lemma \ref{cork1} implies $\rk(T_1)\leq1$. Then,
\begin{align*}
\deg(S_2) & >2\deg(M)-\deg(T_1)-\deg(S_1) \\
& \geq 4\mu_{\max}(V)-\mu_{\max}(V)-\mu_{\max}(V)=2\mu_{\max}(V),
\end{align*}
and this contradicts Lemma \ref{lagrange}. Therefore, we~get that $ \nabla(M)\not\subset M^\perp\otimes\Omega_X(\log D)$.
\end{proof}

Now apply the above results to algebraic curves. For elliptic curves, structure of vector bundles is clear \cite{A57,T93}. We~have the following corollaries.

\begin{cor}
Let $X$ be an elliptic curve. Let $(V,\langle \,, \, \rangle,\nabla)$ be an orthogonal sheaf of rank $5$ together with an orthogonal \lconnection on $X$. Then $(V,\langle \,, \, \rangle,\nabla)$ is not quasi \grsemistable if and only if
\begin{align*}
V&=M\oplus M^\vee\oplus N,
\\
\nabla(M)&\not\subset (M\oplus N)\otimes\Omega_X(\log D),
\end{align*}
where $M$ is a rank two stable bundle of odd degree and $N$ is a line bundle such that $N^{\otimes 2}\cong \mathcal{O}_X$.
\end{cor}

\begin{cor}\label{elliptic6}
Let $X$ be an elliptic curve. Let $(V,\langle \,, \, \rangle,\nabla)$ be an orthogonal sheaf of rank~$6$ together with an orthogonal \lconnection on $X$. Then $(V,\langle \,, \, \rangle,\nabla)$ is not quasi \grsemistable if and only if
\begin{align*}
V&=M\oplus M^\vee\oplus N_1\oplus N_2,
\\
\nabla(M)&\not\subset (M\oplus N_1\oplus N_2)\otimes\Omega_X(\log D),
\end{align*}
where $M$ is a rank two stable bundle of odd degree, and $N_1,N_2$ are non isomorphic line bundles such that $N_i^{\otimes 2}\cong \mathcal{O}_X$ for $i=1,2$.
\end{cor}

For projective line $\mathbb{P}^1_k$, we~have the following corollary. 
\begin{cor}\label{proline5}
Let $(V,\langle \,, \, \rangle,\nabla)$ be an orthogonal bundle on $\mathbb{P}^1_k$ together with an orthogonal \lconnection. If $\rk(V) \leq 6$, then it is quasi \grsemistable.
\end{cor}

\begin{proof}
Note that there is no stable bundle of rank two on $\mathbb{P}^1_k$. This corollary is a direct consequence of Proposition \ref{classfication5} and \ref{classfication6}.
\end{proof}

Based on Corollary \ref{proline5}, we~intend to believe the truth of the following question.

\begin{ques}
Is any orthogonal bundle with an orthogonal \lconnection on the projective line $\PP^1_k$ quasi \grsemistable and hence \grsemistable?
\end{ques}\label{prolineconj}


% End literal retained input author-body.tex

\begin{thebibliography}{99}
\bibitem{AHLH23} J. Alper, D. Halpern-Leistner \& J. Heinloth - “Existence of moduli spaces for algebraic stacks” , Invent. Math. 234 (2023) no. 3, p. 949-1038
\bibitem{A57} M. F. Atiyah - “Vector bundles over an elliptic curve” , Proc. London Math. Soc. (3) 7 (1957), p. 414-452
\bibitem{BMW11} I. Biswas, S. Majumder \& M. L. Wong - “Orthogonal and symplectic parabolic bundles” , J. Geom. Phys. 61 (2011) no. 8, p. 1462-1475
\bibitem{BP03} V. Balaji \& A. J. Parameswaran - “Semistable principal bundles. II. Positive characteristics” , Transform. Groups 8 (2003) no. 1, p. 3-36
\bibitem{CH23} I. Choe \& G. H. Hitching - “Low rank orthogonal bundles and quadric fibrations” , J. Korean Math. Soc. 60 (2023) no. 6, p. 1137-1169
\bibitem{CS21} B. Collier \& A. Sanders - “(G,P)-opers and global Slodowy slices” , Adv. Math. 377 (2021), article ID 107490, 43 pages
\bibitem{Gi73} D. Gieseker - “Stable vector bundles and the Frobenius morphism” , Ann. Sci. École Norm. Sup. (4) 6 (1973), p. 95-101
\bibitem{GS03} T. L. Gómez \& I. Sols - “Stable tensors and moduli space of orthogonal sheaves” , 2003
\bibitem{HL97} D. Huybrechts \& M. Lehn - The geometry of moduli spaces of sheaves , Cambridge Math. Library , Cambridge University Press, Cambridge, 2010
\bibitem{KSZ21} G. Kydonakis, H. Sun \& L. Zhao - “Topological invariants of parabolic $G$-Higgs bundles” , Math. Z. 297 (2021) no. 1-2, p. 585-632
\bibitem{KSZ23} G. Kydonakis, H. Sun \& L. Zhao - “Logahoric Higgs torsors for a complex reductive group” , Math. Ann. 388 (2024) no. 3, p. 3183-3228
\bibitem{L13} A. Langer - “Semistable modules over Lie algebroids in positive characteristic” , Doc. Math. 19 (2014), p. 509-540
\bibitem{Lo17a} T. Lovering - “Integral canonical models for automorphic vector bundles of abelian type” , Algebra Number Theory 11 (2017) no. 8, p. 1837-1890
\bibitem{Lo17b} T. Lovering - “Filtered $F$-crystals on Shimura varieties of abelian type” , 2017
\bibitem{LSZ19} G. Lan, M. Sheng \& K. Zuo - “Semistable Higgs bundles, periodic Higgs bundles and representations of algebraic fundamental groups” , J. Eur. Math. Soc. (JEMS) 21 (2019) no. 10, p. 3053-3112
\bibitem{Ra75} A. Ramanathan - “Stable principal bundles on a compact Riemann surface” , Math. Ann. 213 (1975), p. 129-152
\bibitem{Ra96a} A. Ramanathan - “Moduli for principal bundles over algebraic curves. I” , Proc. Indian Acad. Sci. Math. Sci. 106 (1996) no. 3, p. 301-328
\bibitem{Ra96b} A. Ramanathan - “Moduli for principal bundles over algebraic curves. II” , Proc. Indian Acad. Sci. Math. Sci. 106 (1996) no. 4, p. 421-449
\bibitem{Si92} C. T. Simpson - “Higgs bundles and local systems” , Publ. Math. Inst. Hautes Études Sci. (1992) no. 75, p. 5-95
\bibitem{Si97} C. T. Simpson - “The Hodge filtration on nonabelian cohomology” , in Algebraic geometry (Santa Cruz, 1995) , Proc. Sympos. Pure Math. , vol. 62, Part 2 , American Mathematical Society, Providence, RI, 1997, p. 217-281
\bibitem{S10} C. T. Simpson - “Iterated destabilizing modifications for vector bundles with connection” , in Vector bundles and complex geometry , Contemp. Math. , vol. 522 , American Mathematical Society, Providence, RI, 2010, p. 183-206
\bibitem{T93} L. W. Tu - “Semistable bundles over an elliptic curve” , Adv. Math. 98 (1993) no. 1, p. 1-26
\bibitem{Ya22} M. Yang - “A comparison of generalized opers and $(G,P)$-opers” , Indian J. Pure Appl. Math. 53 (2022) no. 3, p. 760-773
\end{thebibliography}
\end{document}
